\documentclass[]{article}
\usepackage{lmodern}
\usepackage{amssymb,amsmath}
\usepackage{ifxetex,ifluatex}
\usepackage{fixltx2e} % provides \textsubscript
\ifnum 0\ifxetex 1\fi\ifluatex 1\fi=0 % if pdftex
  \usepackage[T1]{fontenc}
  \usepackage[utf8]{inputenc}
\else % if luatex or xelatex
  \ifxetex
    \usepackage{mathspec}
    \usepackage{xltxtra,xunicode}
  \else
    \usepackage{fontspec}
  \fi
  \defaultfontfeatures{Mapping=tex-text,Scale=MatchLowercase}
  \newcommand{\euro}{€}
\fi
% use upquote if available, for straight quotes in verbatim environments
\IfFileExists{upquote.sty}{\usepackage{upquote}}{}
% use microtype if available
\IfFileExists{microtype.sty}{\usepackage{microtype}}{}
\usepackage[margin=1in]{geometry}
\usepackage{color}
\usepackage{fancyvrb}
\newcommand{\VerbBar}{|}
\newcommand{\VERB}{\Verb[commandchars=\\\{\}]}
\DefineVerbatimEnvironment{Highlighting}{Verbatim}{commandchars=\\\{\}}
% Add ',fontsize=\small' for more characters per line
\newenvironment{Shaded}{}{}
\newcommand{\AlertTok}[1]{\textcolor[rgb]{1.00,0.00,0.00}{\textbf{#1}}}
\newcommand{\AnnotationTok}[1]{\textcolor[rgb]{0.38,0.63,0.69}{\textbf{\textit{#1}}}}
\newcommand{\AttributeTok}[1]{\textcolor[rgb]{0.49,0.56,0.16}{#1}}
\newcommand{\BaseNTok}[1]{\textcolor[rgb]{0.25,0.63,0.44}{#1}}
\newcommand{\BuiltInTok}[1]{\textcolor[rgb]{0.00,0.50,0.00}{#1}}
\newcommand{\CharTok}[1]{\textcolor[rgb]{0.25,0.44,0.63}{#1}}
\newcommand{\CommentTok}[1]{\textcolor[rgb]{0.38,0.63,0.69}{\textit{#1}}}
\newcommand{\CommentVarTok}[1]{\textcolor[rgb]{0.38,0.63,0.69}{\textbf{\textit{#1}}}}
\newcommand{\ConstantTok}[1]{\textcolor[rgb]{0.53,0.00,0.00}{#1}}
\newcommand{\ControlFlowTok}[1]{\textcolor[rgb]{0.00,0.44,0.13}{\textbf{#1}}}
\newcommand{\DataTypeTok}[1]{\textcolor[rgb]{0.56,0.13,0.00}{#1}}
\newcommand{\DecValTok}[1]{\textcolor[rgb]{0.25,0.63,0.44}{#1}}
\newcommand{\DocumentationTok}[1]{\textcolor[rgb]{0.73,0.13,0.13}{\textit{#1}}}
\newcommand{\ErrorTok}[1]{\textcolor[rgb]{1.00,0.00,0.00}{\textbf{#1}}}
\newcommand{\ExtensionTok}[1]{#1}
\newcommand{\FloatTok}[1]{\textcolor[rgb]{0.25,0.63,0.44}{#1}}
\newcommand{\FunctionTok}[1]{\textcolor[rgb]{0.02,0.16,0.49}{#1}}
\newcommand{\ImportTok}[1]{\textcolor[rgb]{0.00,0.50,0.00}{\textbf{#1}}}
\newcommand{\InformationTok}[1]{\textcolor[rgb]{0.38,0.63,0.69}{\textbf{\textit{#1}}}}
\newcommand{\KeywordTok}[1]{\textcolor[rgb]{0.00,0.44,0.13}{\textbf{#1}}}
\newcommand{\NormalTok}[1]{#1}
\newcommand{\OperatorTok}[1]{\textcolor[rgb]{0.40,0.40,0.40}{#1}}
\newcommand{\OtherTok}[1]{\textcolor[rgb]{0.00,0.44,0.13}{#1}}
\newcommand{\PreprocessorTok}[1]{\textcolor[rgb]{0.74,0.48,0.00}{#1}}
\newcommand{\RegionMarkerTok}[1]{#1}
\newcommand{\SpecialCharTok}[1]{\textcolor[rgb]{0.25,0.44,0.63}{#1}}
\newcommand{\SpecialStringTok}[1]{\textcolor[rgb]{0.73,0.40,0.53}{#1}}
\newcommand{\StringTok}[1]{\textcolor[rgb]{0.25,0.44,0.63}{#1}}
\newcommand{\VariableTok}[1]{\textcolor[rgb]{0.10,0.09,0.49}{#1}}
\newcommand{\VerbatimStringTok}[1]{\textcolor[rgb]{0.25,0.44,0.63}{#1}}
\newcommand{\WarningTok}[1]{\textcolor[rgb]{0.38,0.63,0.69}{\textbf{\textit{#1}}}}
\ifxetex
  \usepackage[setpagesize=false, % page size defined by xetex
              unicode=false, % unicode breaks when used with xetex
              xetex]{hyperref}
\else
  \usepackage[unicode=true]{hyperref}
\fi
\hypersetup{breaklinks=true,
            bookmarks=true,
            pdfauthor={Justin Le},
            pdftitle={dhall-jit: A Bytecode Compiler and Interpreter for Dhall},
            colorlinks=true,
            citecolor=blue,
            urlcolor=blue,
            linkcolor=magenta,
            pdfborder={0 0 0}}
\urlstyle{same}  % don't use monospace font for urls
% Make links footnotes instead of hotlinks:
\renewcommand{\href}[2]{#2\footnote{\url{#1}}}
\setlength{\parindent}{0pt}
\setlength{\parskip}{6pt plus 2pt minus 1pt}
\setlength{\emergencystretch}{3em}  % prevent overfull lines
\setcounter{secnumdepth}{0}

\title{dhall-jit: A Bytecode Compiler and Interpreter for Dhall}
\author{Justin Le}

\begin{document}
\maketitle

\emph{Originally posted on
\textbf{\href{https://blog.jle.im/entry/dhall-jit-bytecode-compiler-interpreter.html}{in
Code}}.}

\documentclass[]{}
\usepackage{lmodern}
\usepackage{amssymb,amsmath}
\usepackage{ifxetex,ifluatex}
\usepackage{fixltx2e} % provides \textsubscript
\ifnum 0\ifxetex 1\fi\ifluatex 1\fi=0 % if pdftex
  \usepackage[T1]{fontenc}
  \usepackage[utf8]{inputenc}
\else % if luatex or xelatex
  \ifxetex
    \usepackage{mathspec}
    \usepackage{xltxtra,xunicode}
  \else
    \usepackage{fontspec}
  \fi
  \defaultfontfeatures{Mapping=tex-text,Scale=MatchLowercase}
  \newcommand{\euro}{€}
\fi
% use upquote if available, for straight quotes in verbatim environments
\IfFileExists{upquote.sty}{\usepackage{upquote}}{}
% use microtype if available
\IfFileExists{microtype.sty}{\usepackage{microtype}}{}
\usepackage[margin=1in]{geometry}
\ifxetex
  \usepackage[setpagesize=false, % page size defined by xetex
              unicode=false, % unicode breaks when used with xetex
              xetex]{hyperref}
\else
  \usepackage[unicode=true]{hyperref}
\fi
\hypersetup{breaklinks=true,
            bookmarks=true,
            pdfauthor={},
            pdftitle={},
            colorlinks=true,
            citecolor=blue,
            urlcolor=blue,
            linkcolor=magenta,
            pdfborder={0 0 0}}
\urlstyle{same}  % don't use monospace font for urls
% Make links footnotes instead of hotlinks:
\renewcommand{\href}[2]{#2\footnote{\url{#1}}}
\setlength{\parindent}{0pt}
\setlength{\parskip}{6pt plus 2pt minus 1pt}
\setlength{\emergencystretch}{3em}  % prevent overfull lines
\setcounter{secnumdepth}{0}

\textbackslash begin\{document\}

\section{What this is}\label{what-this-is}

A bytecode compiler and multi-backend execution engine for Dhall. Compiles Dhall
expressions to a stack-based bytecode, then executes via a pure Haskell
interpreter or a C interpreter via FFI.

\begin{verbatim}
                Dhall Expr
               /          \
   compileExpr             normalize (dhall library)
             /              \
            v                v
       ProgramBanks        Expr (NF)
            |
     +------+------+
     |             |
     v             v
  Haskell       C interp
  interp         (FFI)
\end{verbatim}

The motivation: Dhall's built-in normalizer produces Dhall \texttt{Expr} values,
which means running a Dhall function requires re-normalizing the full expression
tree on every application. For functions called in hot loops (configuration
functions applied per-element), this is unacceptably slow. The bytecode compiler
pays a one-time cost to compile the function into a flat representation, then
the interpreter can execute repeated applications at near-native speed without
re-traversing syntax trees.

\section{The Bytecode}\label{the-bytecode}

\subsection{ProgramBanks}\label{programbanks}

The compiled representation is a map from integer procedure IDs to instruction
sequences:

\begin{Shaded}
\begin{Highlighting}[]
\KeywordTok{type} \DataTypeTok{ProgramBanks} \OtherTok{=} \DataTypeTok{IntMap}\NormalTok{ (}\DataTypeTok{NESeq}\NormalTok{ (}\DataTypeTok{Instr} \DataTypeTok{Int}\NormalTok{))}
\end{Highlighting}
\end{Shaded}

\texttt{NESeq} (non-empty sequence) because every bank must have at least one
instruction. The \texttt{Int} parameter on \texttt{Instr} is the type of jump
targets -- procedure IDs. Each bank has an implicit RETURN at the end (the final
value left on the stack is the return value).

\begin{itemize}
\tightlist
\item
  Bank 0 is the entry point (MAIN)
\item
  Lambda bodies, if-branches, and let-bindings each get their own bank
\item
  Each bank is a straight-line instruction sequence (no jumps within a bank)
\item
  All branching is done by calling into other banks
\end{itemize}

\subsection{Execution model}\label{execution-model}

Execution state per procedure call:

\begin{itemize}
\tightlist
\item
  A \textbf{data stack} (list of \texttt{Value}) for intermediate values
\item
  A \textbf{read-only environment} (\texttt{Seq\ (Value\ Int)}) -- the captured
  closure variables plus the current argument
\item
  The \textbf{program banks} (immutable code store, shared across all calls)
\end{itemize}

When a procedure finishes executing all its instructions, it pops the single
remaining stack value as its return value.

\subsection{Why no jumps}\label{why-no-jumps}

We can avoid JUMPs completely by treating merges and if/then/else as function
calls into other banks. Every bank is a straight-line sequence with an implicit
RETURN at the end. This means:

\begin{itemize}
\tightlist
\item
  No control-flow analysis needed
\item
  Each bank is independently compilable
\item
  The C interpreter maps trivially: one C function per bank
\item
  Serialization is trivial (just serialize each bank independently)
\end{itemize}

\subsection{Supporting types}\label{supporting-types}

\begin{Shaded}
\begin{Highlighting}[]
\CommentTok{{-}{-} Primitives: the leaf values of the Dhall type system}
\KeywordTok{data} \DataTypeTok{Prim}
    \OtherTok{=} \DataTypeTok{PNat} \OperatorTok{!}\DataTypeTok{Natural}
    \OperatorTok{|} \DataTypeTok{PInt} \OperatorTok{!}\DataTypeTok{Integer}
    \OperatorTok{|} \DataTypeTok{PDouble} \OperatorTok{!}\DataTypeTok{Double}
    \OperatorTok{|} \DataTypeTok{PBool} \OperatorTok{!}\DataTypeTok{Bool}
    \OperatorTok{|} \DataTypeTok{PText} \OperatorTok{!}\DataTypeTok{Text}
    \OperatorTok{|} \DataTypeTok{PBytes} \OperatorTok{!}\DataTypeTok{ByteString}
    \OperatorTok{|} \DataTypeTok{PDate} \OperatorTok{!}\DataTypeTok{Day}
    \OperatorTok{|} \DataTypeTok{PTime} \OperatorTok{!}\DataTypeTok{TimeOfDay}
    \OperatorTok{|} \DataTypeTok{PTimeZone} \OperatorTok{!}\DataTypeTok{TimeZone}
    \OperatorTok{|} \DataTypeTok{PAssert}
    \OperatorTok{|} \DataTypeTok{PType}
    \CommentTok{{-}{-} PType is a dummy value used to fill stack slots for type arguments.}
    \CommentTok{{-}{-} Dhall has dependent types at the kind level: (\textbackslash{}(a : Type) {-}\textgreater{} \textbackslash{}(x : a) {-}\textgreater{} x)}
    \CommentTok{{-}{-} needs *something* on the stack for the Type argument so de Bruijn indices}
    \CommentTok{{-}{-} line up. PType occupies that slot and is never inspected.}

\CommentTok{{-}{-} Path component for WITH instruction (nested record update)}
\KeywordTok{data} \DataTypeTok{PathComponent}
    \OtherTok{=} \DataTypeTok{PathLabel} \OperatorTok{!}\DataTypeTok{Word16}  \CommentTok{{-}{-} field index}
    \OperatorTok{|} \DataTypeTok{PathQuestion}       \CommentTok{{-}{-} Optional\textquotesingle{}s Some case}

\CommentTok{{-}{-} Arity of a union alternative (used by MERGE\_UNION)}
\KeywordTok{data} \DataTypeTok{UnionArity} \OtherTok{=} \DataTypeTok{Nullary} \OperatorTok{|} \DataTypeTok{Unary}

\CommentTok{{-}{-} How to merge a field from two records}
\KeywordTok{data} \DataTypeTok{RecordMergeOp}\NormalTok{ a}
    \OtherTok{=} \DataTypeTok{FromLeft} \OperatorTok{!}\NormalTok{a    }\CommentTok{{-}{-} take from left operand at index a}
    \OperatorTok{|} \DataTypeTok{FromRight} \OperatorTok{!}\NormalTok{a   }\CommentTok{{-}{-} take from right operand at index a}

\CommentTok{{-}{-} Builtin operations and functions}
\KeywordTok{data} \DataTypeTok{DhallBuiltin}
    \OtherTok{=} \DataTypeTok{ADD\_NAT} \OperatorTok{|} \DataTypeTok{MUL\_NAT} \OperatorTok{|} \DataTypeTok{TEXT\_APPEND} \OperatorTok{|} \DataTypeTok{LIST\_APPEND}
    \OperatorTok{|} \DataTypeTok{AND\_BOOL} \OperatorTok{|} \DataTypeTok{OR\_BOOL} \OperatorTok{|} \DataTypeTok{EQ\_BOOL} \OperatorTok{|} \DataTypeTok{NE\_BOOL}
    \OperatorTok{|} \DataTypeTok{NATURAL\_FOLD} \OperatorTok{|} \DataTypeTok{NATURAL\_BUILD}
    \OperatorTok{|} \DataTypeTok{NATURAL\_IS\_ZERO} \OperatorTok{|} \DataTypeTok{NATURAL\_EVEN} \OperatorTok{|} \DataTypeTok{NATURAL\_ODD}
    \OperatorTok{|} \DataTypeTok{NATURAL\_TO\_INTEGER} \OperatorTok{|} \DataTypeTok{NATURAL\_SHOW} \OperatorTok{|} \DataTypeTok{NATURAL\_SUBTRACT}
    \OperatorTok{|} \DataTypeTok{INTEGER\_CLAMP} \OperatorTok{|} \DataTypeTok{INTEGER\_NEGATE} \OperatorTok{|} \DataTypeTok{INTEGER\_SHOW} \OperatorTok{|} \DataTypeTok{INTEGER\_TO\_DOUBLE}
    \OperatorTok{|} \DataTypeTok{DOUBLE\_SHOW} \OperatorTok{|} \DataTypeTok{TEXT\_SHOW} \OperatorTok{|} \DataTypeTok{TEXT\_REPLACE}
    \OperatorTok{|} \DataTypeTok{DATE\_SHOW} \OperatorTok{|} \DataTypeTok{TIME\_SHOW} \OperatorTok{|} \DataTypeTok{TIMEZONE\_SHOW}
    \OperatorTok{|} \DataTypeTok{LIST\_BUILD} \OperatorTok{|} \DataTypeTok{LIST\_FOLD}
    \OperatorTok{|} \DataTypeTok{LIST\_LENGTH} \OperatorTok{|} \DataTypeTok{LIST\_HEAD} \OperatorTok{|} \DataTypeTok{LIST\_LAST} \OperatorTok{|} \DataTypeTok{LIST\_INDEXED} \OperatorTok{|} \DataTypeTok{LIST\_REVERSE}
    \OperatorTok{|} \DataTypeTok{NONE}

\CommentTok{{-}{-} Functions provided by the runtime to user code (for fold/build)}
\KeywordTok{data} \DataTypeTok{ProvidedFn} \OtherTok{=} \DataTypeTok{NaturalSucc} \OperatorTok{|} \DataTypeTok{ListCons}
\end{Highlighting}
\end{Shaded}

\subsection{The instruction set}\label{the-instruction-set}

\begin{Shaded}
\begin{Highlighting}[]
\KeywordTok{data} \DataTypeTok{Instr}\NormalTok{ i}
    \OtherTok{=} \DataTypeTok{LOAD\_LOCAL} \OperatorTok{!}\DataTypeTok{Word8}
    \CommentTok{{-}{-} Push env[i] onto the stack.}
    \CommentTok{{-}{-} Index 0 is the current argument (innermost lambda parameter).}
    \CommentTok{{-}{-} Index 1+ are captured variables from the enclosing scope.}

    \OperatorTok{|} \DataTypeTok{MAKE\_LIST} \OperatorTok{!}\DataTypeTok{Word16}
    \CommentTok{{-}{-} Pop N values from the stack (topmost = last element),}
    \CommentTok{{-}{-} construct a VList. N=0 is valid (empty list).}

    \OperatorTok{|} \DataTypeTok{MAKE\_RECORD} \OperatorTok{!}\DataTypeTok{Word16}
    \CommentTok{{-}{-} Pop N values from the stack (topmost = last field),}
    \CommentTok{{-}{-} construct a VRecord. Fields are in alphabetically{-}sorted order}
    \CommentTok{{-}{-} (matching Dhall\textquotesingle{}s canonical field ordering).}

    \OperatorTok{|} \DataTypeTok{LOAD\_FIELD} \OperatorTok{!}\DataTypeTok{Word16}
    \CommentTok{{-}{-} Pop a record, push the field at the given index.}

    \OperatorTok{|} \DataTypeTok{MAKE\_UNION} \OperatorTok{!}\DataTypeTok{Word8}
    \CommentTok{{-}{-} Pop a value, wrap it as VUnion with the given tag.}

    \OperatorTok{|} \DataTypeTok{MAKE\_NULLARY\_UNION} \OperatorTok{!}\DataTypeTok{Word8}
    \CommentTok{{-}{-} Push VUnion tag Nothing (no pop {-}{-} the constructor has no payload).}

    \OperatorTok{|} \DataTypeTok{MERGE\_UNION} \OperatorTok{!}\NormalTok{(}\DataTypeTok{Vector} \DataTypeTok{UnionArity}\NormalTok{)}
    \CommentTok{{-}{-} Pop a union value, then pop a handler record.}
    \CommentTok{{-}{-} The Vector gives the arity of each alternative (indexed by tag).}
    \CommentTok{{-}{-} If Nullary: push the handler directly.}
    \CommentTok{{-}{-} If Unary: CALL the handler with the union\textquotesingle{}s payload.}

    \OperatorTok{|} \DataTypeTok{MERGE\_IF} \OperatorTok{!}\NormalTok{i }\OperatorTok{!}\NormalTok{i}
    \CommentTok{{-}{-} Pop a Bool. If True, execute bank i1; if False, execute bank i2.}
    \CommentTok{{-}{-} Both banks inherit the current environment.}

    \OperatorTok{|} \DataTypeTok{PUSH\_PRIM} \OperatorTok{!}\DataTypeTok{Prim}
    \CommentTok{{-}{-} Push a literal value.}

    \OperatorTok{|} \DataTypeTok{PUSH\_BUILTIN} \OperatorTok{!}\DataTypeTok{DhallBuiltin}
    \CommentTok{{-}{-} Push VBuiltin with zero captured arguments.}

    \OperatorTok{|} \DataTypeTok{BINOP} \OperatorTok{!}\DataTypeTok{DhallBuiltin}
    \CommentTok{{-}{-} Pop two values [left, right] (right on top), apply the binary}
    \CommentTok{{-}{-} operation, push the result. Only used for: ADD\_NAT, MUL\_NAT,}
    \CommentTok{{-}{-} TEXT\_APPEND, LIST\_APPEND, AND\_BOOL, OR\_BOOL, EQ\_BOOL, NE\_BOOL.}

    \OperatorTok{|} \DataTypeTok{CLOSURE}\NormalTok{ i}
    \CommentTok{{-}{-} Capture the entire current environment, push VLam i env.}

    \OperatorTok{|} \DataTypeTok{CALL}
    \CommentTok{{-}{-} Pop arg (top), pop fn (second).}
    \CommentTok{{-}{-} If fn is VLam: execute bank with env = arg \textless{}| capturedEnv.}
    \CommentTok{{-}{-} If fn is VBuiltin: accumulate arg; fire if saturated.}
    \CommentTok{{-}{-} If fn is VProvided: accumulate arg; fire if saturated.}

    \OperatorTok{|} \DataTypeTok{SHOW\_CONSTRUCTOR} \OperatorTok{!}\NormalTok{(}\DataTypeTok{Vector} \DataTypeTok{Text}\NormalTok{)}
    \CommentTok{{-}{-} Pop a union, push the constructor name (as PText) by indexing}
    \CommentTok{{-}{-} into the vector with the union\textquotesingle{}s tag.}

    \OperatorTok{|} \DataTypeTok{TOMAP} \OperatorTok{!}\NormalTok{(}\DataTypeTok{Vector} \DataTypeTok{Text}\NormalTok{)}
    \CommentTok{{-}{-} Pop a record, produce a list of \{mapKey: Text, mapValue: a\}}
    \CommentTok{{-}{-} records. Vector contains field names in sorted order.}

    \OperatorTok{|} \DataTypeTok{WITH} \OperatorTok{!}\NormalTok{(}\DataTypeTok{NESeq} \DataTypeTok{PathComponent}\NormalTok{)}
    \CommentTok{{-}{-} Pop new value (top), pop base value (second).}
    \CommentTok{{-}{-} Traverse base along the path, replacing the final target with}
    \CommentTok{{-}{-} the new value. PathLabel descends into record fields; PathQuestion}
    \CommentTok{{-}{-} descends into Some.}

    \OperatorTok{|} \DataTypeTok{PROJECT} \OperatorTok{!}\NormalTok{(}\DataTypeTok{Vector} \DataTypeTok{Word16}\NormalTok{)}
    \CommentTok{{-}{-} Pop a record, push a new record containing only the fields at}
    \CommentTok{{-}{-} the given indices (in sorted order).}

    \OperatorTok{|} \DataTypeTok{COMBINE} \OperatorTok{!}\NormalTok{(}\DataTypeTok{Vector}\NormalTok{ (}\DataTypeTok{RecordMergeOp}\NormalTok{ (}\DataTypeTok{NESeq} \DataTypeTok{Word16}\NormalTok{)))}
    \CommentTok{{-}{-} Pop right record (top), pop left record (second).}
    \CommentTok{{-}{-} Each op says: for this output field, take from left or right at}
    \CommentTok{{-}{-} the given path. NESeq Word16 is a path of indices for recursive}
    \CommentTok{{-}{-} merging into nested records.}

    \OperatorTok{|} \DataTypeTok{PREFER} \OperatorTok{!}\NormalTok{(}\DataTypeTok{Vector}\NormalTok{ (}\DataTypeTok{RecordMergeOp} \DataTypeTok{Word16}\NormalTok{))}
    \CommentTok{{-}{-} Pop right record (top), pop left record (second).}
    \CommentTok{{-}{-} Each op says: take field from left at idx, or from right at idx.}
    \CommentTok{{-}{-} Right wins on collision (Dhall\textquotesingle{}s // semantics).}
\end{Highlighting}
\end{Shaded}

\subsection{Runtime values}\label{runtime-values}

\begin{Shaded}
\begin{Highlighting}[]
\KeywordTok{data} \DataTypeTok{Value}\NormalTok{ i}
    \OtherTok{=} \DataTypeTok{VPrim} \OperatorTok{!}\DataTypeTok{Prim}
    \OperatorTok{|} \DataTypeTok{VRecord} \OperatorTok{!}\NormalTok{(}\DataTypeTok{Vector}\NormalTok{ (}\DataTypeTok{Value}\NormalTok{ i))}
    \OperatorTok{|} \DataTypeTok{VList} \OperatorTok{!}\NormalTok{(}\DataTypeTok{Vector}\NormalTok{ (}\DataTypeTok{Value}\NormalTok{ i))}
    \OperatorTok{|} \DataTypeTok{VUnion} \OperatorTok{!}\DataTypeTok{Word8} \OperatorTok{!}\NormalTok{(}\DataTypeTok{Maybe}\NormalTok{ (}\DataTypeTok{Value}\NormalTok{ i))}
    \OperatorTok{|} \DataTypeTok{VLam} \OperatorTok{!}\NormalTok{i (}\DataTypeTok{Seq}\NormalTok{ (}\DataTypeTok{Value}\NormalTok{ i))}
    \OperatorTok{|} \DataTypeTok{VBuiltin} \OperatorTok{!}\DataTypeTok{DhallBuiltin}\NormalTok{ [}\DataTypeTok{Value}\NormalTok{ i]}
    \OperatorTok{|} \DataTypeTok{VProvided} \OperatorTok{!}\DataTypeTok{ProvidedFn}\NormalTok{ [}\DataTypeTok{Value}\NormalTok{ i]}
\end{Highlighting}
\end{Shaded}

Key difference from a normalized Dhall Expr: no \texttt{App} (everything is
fully-applied), and lambdas are jumps with closures rather than syntax trees.

\texttt{VBuiltin} and \texttt{VProvided} accumulate arguments in a list. When a
\texttt{CALL} provides an argument to a \texttt{VBuiltin}, it appends the arg.
Once the arg count reaches the builtin's arity, the builtin executes and
produces a result value. This means partially-applied builtins are first-class
values -- \texttt{Natural/add\ 3} is
\texttt{VBuiltin\ ADD\_NAT\ {[}VPrim\ (PNat\ 3){]}} sitting on the stack,
waiting for one more arg.

\texttt{VProvided} exists for the two functions that \texttt{Natural/fold} and
\texttt{List/build} pass to user code: \texttt{NaturalSucc} (increments a
Natural) and \texttt{ListCons} (prepends to a list). These are
runtime-synthesized functions, not user-written lambdas.

\section{The Compiler (Compile.hs)}\label{the-compiler-compile.hs}

The compiler works in a single pass over a type-checked, alpha-normalized Dhall
expression. It emits instructions into the current bank, allocating new banks
for sub-expressions that need their own scope.

\subsection{Compiler state}\label{compiler-state}

\begin{Shaded}
\begin{Highlighting}[]
\KeywordTok{newtype} \DataTypeTok{CompileState} \OtherTok{=} \DataTypeTok{CS}\NormalTok{ \{}\OtherTok{csBanks ::} \DataTypeTok{ProgramBanks}\NormalTok{\}}

\KeywordTok{data} \DataTypeTok{CompileContext}\NormalTok{ s }\OtherTok{=} \DataTypeTok{CC}
\NormalTok{    \{}\OtherTok{ ccDhallCtx ::} \OperatorTok{!}\NormalTok{(}\DataTypeTok{Context}\NormalTok{ (}\DataTypeTok{FunctionBinding}\NormalTok{ s }\DataTypeTok{Void}\NormalTok{))}
    \CommentTok{{-}{-} \^{} Dhall typing context: tracks what variables are in scope and their types.}
    \CommentTok{{-}{-}   Used for type{-}directed compilation (field lookups, union arities).}
\NormalTok{    ,}\OtherTok{ ccCurrBank ::} \OperatorTok{!}\DataTypeTok{Int}
    \CommentTok{{-}{-} \^{} Which bank we are currently emitting instructions into.}
\NormalTok{    \}}

\KeywordTok{type} \DataTypeTok{Compile}\NormalTok{ s }\OtherTok{=} \DataTypeTok{ReaderT}\NormalTok{ (}\DataTypeTok{CompileContext}\NormalTok{ s) (}\DataTypeTok{StateT} \DataTypeTok{CompileState}\NormalTok{ (}\DataTypeTok{Except}\NormalTok{ (}\DataTypeTok{CompileError}\NormalTok{ s)))}
\end{Highlighting}
\end{Shaded}

\subsection{Bank allocation}\label{bank-allocation}

\begin{Shaded}
\begin{Highlighting}[]
\CommentTok{{-}{-} Allocate a new bank ID (one past the current max)}
\OtherTok{allocBank ::} \DataTypeTok{Compile}\NormalTok{ s }\DataTypeTok{Int}

\CommentTok{{-}{-} Run an action emitting into a fresh bank, return the bank ID}
\OtherTok{withNewBank\_ ::} \DataTypeTok{Compile}\NormalTok{ s a }\OtherTok{{-}\textgreater{}} \DataTypeTok{Compile}\NormalTok{ s }\DataTypeTok{Int}
\end{Highlighting}
\end{Shaded}

\subsection{How each Dhall construct
compiles}\label{how-each-dhall-construct-compiles}

\textbf{Variables} (\texttt{Var\ (V\ \_\ idx)}):

\begin{verbatim}
emit $ LOAD_LOCAL (fromIntegral idx)
\end{verbatim}

De Bruijn index directly becomes the environment offset.

\textbf{Lambdas} (\texttt{Lam\ \_\ binding\ body}):

\begin{verbatim}
bodyBank <- withNewBank_ $ compileExpr_ body  -- body goes into its own bank
emit $ CLOSURE bodyBank                        -- capture current env
\end{verbatim}

The lambda body is compiled into a separate bank. At runtime,
\texttt{CLOSURE\ bodyBank} captures the current environment and pushes a
\texttt{VLam\ bodyBank\ env}. When this closure is later \texttt{CALL}ed, the
interpreter enters \texttt{bodyBank} with
\texttt{arg\ \textless{}\textbar{}\ env} as the new environment.

\textbf{Application} (\texttt{App\ f\ x}):

\begin{verbatim}
compileExpr_ f          -- push the function
compileExpr_ x          -- push the argument (or PType if x is a Type)
emit CALL               -- invoke
\end{verbatim}

Type arguments are detected at compile time: if the argument's type is
\texttt{Const\ Type/Kind/Sort}, we push \texttt{PType} instead of compiling the
actual type expression (it will never be inspected at runtime).

\textbf{Let bindings} (\texttt{Let\ binding\ body}):

\begin{verbatim}
compileExpr_ (value binding)   -- push the bound value
bodyBank <- withNewBank_ $ compileExpr_ body
emit $ CLOSURE bodyBank        -- body as a lambda
emit CALL                      -- immediately apply
\end{verbatim}

\texttt{let\ x\ =\ e1\ in\ e2} is semantically
\texttt{(\textbackslash{}x\ -\textgreater{}\ e2)\ e1}. The compiler emits it
literally that way. The body bank gets the bound value as env{[}0{]}.

\textbf{If/then/else} (\texttt{BoolIf\ b\ t\ f}):

\begin{verbatim}
compileExpr_ b
tBank <- withNewBank_ $ compileExpr_ t
fBank <- withNewBank_ $ compileExpr_ f
emit $ MERGE_IF tBank fBank
\end{verbatim}

Both branches are compiled into separate banks. At runtime, the interpreter pops
the Bool and calls into the chosen bank (inheriting the current env).

\textbf{Record literals} (\texttt{RecordLit\ fields}):

\begin{verbatim}
traverse_ (compileExpr_ . recordFieldValue . snd) (DM.toAscList fields)
emit $ MAKE_RECORD (fromIntegral (DM.size fields))
\end{verbatim}

Fields are compiled in alphabetical order (Dhall's canonical ordering). The
record is positional from this point on.

\textbf{Field access} (\texttt{Field\ record\ fieldName}):

\begin{verbatim}
-- At compile time: look up the field's index in the record type
recordType <- typeOf record
idx <- fieldIndex recordType fieldName
compileExpr_ record
emit $ LOAD_FIELD (fromIntegral idx)
\end{verbatim}

The type checker tells us what position the field occupies. No string comparison
at runtime.

\textbf{Union constructors} (\texttt{Field\ (Union\ alts)\ name}):

For a nullary constructor:

\begin{verbatim}
emit $ MAKE_NULLARY_UNION (fromIntegral idx)
\end{verbatim}

For a unary constructor (takes a payload):

\begin{verbatim}
lambdaBank <- withNewBank_ $ do
    emit $ LOAD_LOCAL 0          -- the payload
    emit $ MAKE_UNION (fromIntegral idx)
emit $ CLOSURE lambdaBank        -- push the constructor as a function
\end{verbatim}

A unary union constructor compiles to a lambda that wraps its argument.

\textbf{Merge} (\texttt{Merge\ handlers\ union\ \_}):

\begin{verbatim}
unionAlts <- typeOf union  -- get the union type to know arities
compileExpr_ handlers
compileExpr_ union
emit $ MERGE_UNION $ V.fromList $ altArity <$> DM.toAscList unionAlts
\end{verbatim}

The arity vector tells the interpreter whether each alternative is Nullary (just
return the handler) or Unary (call the handler with the payload).

\textbf{Prefer} (\texttt{Prefer\ \_\ \_\ l\ r}):

\begin{verbatim}
lFields <- getRecordFields l
rFields <- getRecordFields r
compileExpr_ l
compileExpr_ r
emit $ PREFER $ V.fromList $ buildPreferOps lFields rFields
\end{verbatim}

\texttt{buildPreferOps} produces a vector of
\texttt{FromLeft\ idx\ \textbar{}\ FromRight\ idx} for each output field. On
collision, right wins. The interpreter just iterates the vector and indexes into
the appropriate operand.

\textbf{Combine} (\texttt{Combine\ \_\ \_\ l\ r}):

\begin{verbatim}
lFields <- getRecordFields l
rFields <- getRecordFields r
compileExpr_ l
compileExpr_ r
emit $ COMBINE $ V.fromList $ buildCombineOps lFields rFields
\end{verbatim}

Like PREFER, but recursive. When both sides have a record at the same field,
\texttt{buildCombineOps} encodes the full path as a \texttt{NESeq\ Word16}. The
interpreter follows the path to recursively merge nested records.

\textbf{Text interpolation} (\texttt{TextLit\ (Chunks\ pairs\ suffix)}):

\begin{verbatim}
-- "Hello, ${name}!" compiles as:
-- PUSH_PRIM "Hello, "
-- [compile name]
-- BINOP TEXT_APPEND
-- PUSH_PRIM "!"
-- BINOP TEXT_APPEND
\end{verbatim}

A chain of TEXT\_APPEND operations. Each interpolated chunk pushes a prefix
string, compiles the interpolation, and appends.

\textbf{Builtins} (e.g.~\texttt{NaturalFold}, \texttt{ListBuild}, etc.):

\begin{verbatim}
emit $ PUSH_BUILTIN NATURAL_FOLD
\end{verbatim}

Builtins are values. They sit on the stack as \texttt{VBuiltin\ name\ {[}{]}}
until enough arguments accumulate via \texttt{CALL}.

\textbf{Binary operators} (\texttt{NaturalPlus\ l\ r}, \texttt{BoolAnd\ l\ r},
etc.):

\begin{verbatim}
compileExpr_ l
compileExpr_ r
emit $ BINOP ADD_NAT   -- (or appropriate op)
\end{verbatim}

\section{The Haskell Interpreter
(Interpret.hs)}\label{the-haskell-interpreter-interpret.hs}

\subsection{Execution monad}\label{execution-monad}

\begin{Shaded}
\begin{Highlighting}[]
\KeywordTok{data} \DataTypeTok{ExecContext} \OtherTok{=} \DataTypeTok{ExecContext}
\NormalTok{    \{}\OtherTok{ ecBanks ::} \OperatorTok{!}\DataTypeTok{ProgramBanks}
\NormalTok{    ,}\OtherTok{ ecEnv   ::} \OperatorTok{!}\NormalTok{(}\DataTypeTok{Seq}\NormalTok{ (}\DataTypeTok{Value} \DataTypeTok{Int}\NormalTok{))}
\NormalTok{    \}}

\KeywordTok{type} \DataTypeTok{Exec}\NormalTok{ s }\OtherTok{=} \DataTypeTok{ReaderT} \DataTypeTok{ExecContext}\NormalTok{ (}\DataTypeTok{Except}\NormalTok{ (}\DataTypeTok{ExecError}\NormalTok{ s))}
\end{Highlighting}
\end{Shaded}

\subsection{Core execution loop}\label{core-execution-loop}

\begin{Shaded}
\begin{Highlighting}[]
\OtherTok{executeProcedure ::} \DataTypeTok{Int} \OtherTok{{-}\textgreater{}} \DataTypeTok{Seq}\NormalTok{ (}\DataTypeTok{Value} \DataTypeTok{Int}\NormalTok{) }\OtherTok{{-}\textgreater{}} \DataTypeTok{Exec}\NormalTok{ s (}\DataTypeTok{Value} \DataTypeTok{Int}\NormalTok{)}
\NormalTok{executeProcedure procId newEnv }\OtherTok{=} \KeywordTok{do}
\NormalTok{    instrs }\OtherTok{\textless{}{-}}\NormalTok{ lookupBank procId}
\NormalTok{    local (\textbackslash{}c }\OtherTok{{-}\textgreater{}}\NormalTok{ c\{ecEnv }\OtherTok{=}\NormalTok{ newEnv\}) }\OperatorTok{$}
\NormalTok{        evalStateT (traverse\_ executeInstr (toList instrs) }\OperatorTok{\textgreater{}\textgreater{}}\NormalTok{ pop) []}
\end{Highlighting}
\end{Shaded}

Enter a bank, set the environment, execute all instructions (they push/pop from
the stack), then pop the final value as the result.

\subsection{Instruction execution (key
cases)}\label{instruction-execution-key-cases}

\textbf{CALL} -- the heart of the interpreter:

\begin{Shaded}
\begin{Highlighting}[]
\DataTypeTok{CALL} \OtherTok{{-}\textgreater{}} \KeywordTok{do}
\NormalTok{    arg }\OtherTok{\textless{}{-}}\NormalTok{ pop}
\NormalTok{    fn }\OtherTok{\textless{}{-}}\NormalTok{ pop}
\NormalTok{    push }\OperatorTok{=\textless{}\textless{}} \KeywordTok{case}\NormalTok{ fn }\KeywordTok{of}
        \DataTypeTok{VLam}\NormalTok{ procId closedEnv }\OtherTok{{-}\textgreater{}}
\NormalTok{            lift }\OperatorTok{$}\NormalTok{ executeProcedure procId (arg }\OperatorTok{\textless{}|}\NormalTok{ closedEnv)}
        \DataTypeTok{VBuiltin}\NormalTok{ builtin capturedArgs }\OtherTok{{-}\textgreater{}}
\NormalTok{            lift }\OperatorTok{$}\NormalTok{ applyBuiltin builtin (arg }\OperatorTok{:|}\NormalTok{ capturedArgs)}
        \DataTypeTok{VProvided}\NormalTok{ provFn capturedArgs }\OtherTok{{-}\textgreater{}}
\NormalTok{            lift }\OperatorTok{$}\NormalTok{ executeProvidedFn provFn (arg }\OperatorTok{:|}\NormalTok{ capturedArgs)}
\NormalTok{        \_ }\OtherTok{{-}\textgreater{}}\NormalTok{ throwError }\OperatorTok{$} \DataTypeTok{ExpectedFunction}\NormalTok{ fn}
\end{Highlighting}
\end{Shaded}

For \texttt{VLam}: recursively execute the target bank with \texttt{arg}
prepended to the captured env.

For \texttt{VBuiltin}: add the arg to the list; if now saturated, fire the
builtin.

\textbf{MERGE\_UNION}:

\begin{Shaded}
\begin{Highlighting}[]
\DataTypeTok{MERGE\_UNION}\NormalTok{ arities }\OtherTok{{-}\textgreater{}} \KeywordTok{do}
\NormalTok{    union }\OtherTok{\textless{}{-}}\NormalTok{ pop}
\NormalTok{    handlers }\OtherTok{\textless{}{-}}\NormalTok{ pop}
    \KeywordTok{case}\NormalTok{ (union, handlers) }\KeywordTok{of}
\NormalTok{        (}\DataTypeTok{VUnion}\NormalTok{ tag payload, }\DataTypeTok{VRecord}\NormalTok{ handlerFields) }\OtherTok{{-}\textgreater{}}
            \KeywordTok{case}\NormalTok{ (arities }\OperatorTok{V.!?} \FunctionTok{fromIntegral}\NormalTok{ tag, handlerFields }\OperatorTok{V.!?} \FunctionTok{fromIntegral}\NormalTok{ tag, payload) }\KeywordTok{of}
\NormalTok{                (}\DataTypeTok{Just} \DataTypeTok{Nullary}\NormalTok{, }\DataTypeTok{Just}\NormalTok{ handler, }\DataTypeTok{Nothing}\NormalTok{) }\OtherTok{{-}\textgreater{}}\NormalTok{ push handler}
\NormalTok{                (}\DataTypeTok{Just} \DataTypeTok{Unary}\NormalTok{, }\DataTypeTok{Just}\NormalTok{ handler, }\DataTypeTok{Just}\NormalTok{ payloadVal) }\OtherTok{{-}\textgreater{}} \KeywordTok{do}
                    \CommentTok{{-}{-} Apply handler to payload}
\NormalTok{                    push handler}
\NormalTok{                    push payloadVal}
\NormalTok{                    executeInstr }\DataTypeTok{CALL}
\end{Highlighting}
\end{Shaded}

\textbf{MERGE\_IF}:

\begin{Shaded}
\begin{Highlighting}[]
\DataTypeTok{MERGE\_IF}\NormalTok{ tProc fProc }\OtherTok{{-}\textgreater{}} \KeywordTok{do}
\NormalTok{    cond }\OtherTok{\textless{}{-}}\NormalTok{ pop}
\NormalTok{    push }\OperatorTok{=\textless{}\textless{}} \KeywordTok{case}\NormalTok{ cond }\KeywordTok{of}
        \DataTypeTok{VPrim}\NormalTok{ (}\DataTypeTok{PBool} \DataTypeTok{True}\NormalTok{)  }\OtherTok{{-}\textgreater{}}\NormalTok{ lift }\OperatorTok{$}\NormalTok{ executeProcedure tProc env}
        \DataTypeTok{VPrim}\NormalTok{ (}\DataTypeTok{PBool} \DataTypeTok{False}\NormalTok{) }\OtherTok{{-}\textgreater{}}\NormalTok{ lift }\OperatorTok{$}\NormalTok{ executeProcedure fProc env}
\end{Highlighting}
\end{Shaded}

\subsection{Builtin application}\label{builtin-application}

Builtins fire when all arguments have arrived. Key cases:

\begin{Shaded}
\begin{Highlighting}[]
\OtherTok{applyBuiltin ::} \DataTypeTok{DhallBuiltin} \OtherTok{{-}\textgreater{}} \DataTypeTok{NonEmpty}\NormalTok{ (}\DataTypeTok{Value} \DataTypeTok{Int}\NormalTok{) }\OtherTok{{-}\textgreater{}} \DataTypeTok{Exec}\NormalTok{ s (}\DataTypeTok{Value} \DataTypeTok{Int}\NormalTok{)}
\NormalTok{applyBuiltin builtin allArgs }\OtherTok{=} \KeywordTok{case}\NormalTok{ (builtin, allArgs) }\KeywordTok{of}
    \CommentTok{{-}{-} Arity 1}
\NormalTok{    (}\DataTypeTok{NATURAL\_IS\_ZERO}\NormalTok{, }\DataTypeTok{VPrim}\NormalTok{ (}\DataTypeTok{PNat}\NormalTok{ n) }\OperatorTok{:|}\NormalTok{ []) }\OtherTok{{-}\textgreater{}}
        \FunctionTok{pure} \OperatorTok{$} \DataTypeTok{VPrim} \OperatorTok{$} \DataTypeTok{PBool}\NormalTok{ (n }\OperatorTok{==} \DecValTok{0}\NormalTok{)}
\NormalTok{    (}\DataTypeTok{NATURAL\_TO\_INTEGER}\NormalTok{, }\DataTypeTok{VPrim}\NormalTok{ (}\DataTypeTok{PNat}\NormalTok{ n) }\OperatorTok{:|}\NormalTok{ []) }\OtherTok{{-}\textgreater{}}
        \FunctionTok{pure} \OperatorTok{$} \DataTypeTok{VPrim} \OperatorTok{$} \DataTypeTok{PInt}\NormalTok{ (}\FunctionTok{fromIntegral}\NormalTok{ n)}

    \CommentTok{{-}{-} Arity 2 (note: args accumulate in order, so last applied is head)}
\NormalTok{    (}\DataTypeTok{NATURAL\_SUBTRACT}\NormalTok{, }\DataTypeTok{VPrim}\NormalTok{ (}\DataTypeTok{PNat}\NormalTok{ y) }\OperatorTok{:|}\NormalTok{ [}\DataTypeTok{VPrim}\NormalTok{ (}\DataTypeTok{PNat}\NormalTok{ x)]) }\OtherTok{{-}\textgreater{}}
        \FunctionTok{pure} \OperatorTok{$} \DataTypeTok{VPrim} \OperatorTok{$} \DataTypeTok{PNat}\NormalTok{ (}\KeywordTok{if}\NormalTok{ y }\OperatorTok{\textgreater{}=}\NormalTok{ x }\KeywordTok{then}\NormalTok{ y }\OperatorTok{{-}}\NormalTok{ x }\KeywordTok{else} \DecValTok{0}\NormalTok{)}
\NormalTok{    (}\DataTypeTok{LIST\_LENGTH}\NormalTok{, }\DataTypeTok{VList}\NormalTok{ xs }\OperatorTok{:|}\NormalTok{ [\_]) }\OtherTok{{-}\textgreater{}}
        \FunctionTok{pure} \OperatorTok{$} \DataTypeTok{VPrim} \OperatorTok{$} \DataTypeTok{PNat}\NormalTok{ (}\FunctionTok{fromIntegral}\NormalTok{ (V.length xs))}
\NormalTok{    (}\DataTypeTok{LIST\_HEAD}\NormalTok{, }\DataTypeTok{VList}\NormalTok{ xs }\OperatorTok{:|}\NormalTok{ [\_]) }\OtherTok{{-}\textgreater{}} \KeywordTok{case}\NormalTok{ V.uncons xs }\KeywordTok{of}
        \DataTypeTok{Just}\NormalTok{ (x, \_) }\OtherTok{{-}\textgreater{}} \FunctionTok{pure} \OperatorTok{$} \DataTypeTok{VUnion} \DecValTok{1}\NormalTok{ (}\DataTypeTok{Just}\NormalTok{ x)   }\CommentTok{{-}{-} Some x}
        \DataTypeTok{Nothing}     \OtherTok{{-}\textgreater{}} \FunctionTok{pure} \OperatorTok{$} \DataTypeTok{VUnion} \DecValTok{0} \DataTypeTok{Nothing}     \CommentTok{{-}{-} None}

    \CommentTok{{-}{-} Arity 3}
\NormalTok{    (}\DataTypeTok{TEXT\_REPLACE}\NormalTok{, }\DataTypeTok{VPrim}\NormalTok{ (}\DataTypeTok{PText}\NormalTok{ haystack) }\OperatorTok{:|}\NormalTok{ [}\DataTypeTok{VPrim}\NormalTok{ (}\DataTypeTok{PText}\NormalTok{ replacement), }\DataTypeTok{VPrim}\NormalTok{ (}\DataTypeTok{PText}\NormalTok{ needle)]) }\OtherTok{{-}\textgreater{}}
        \FunctionTok{pure} \OperatorTok{$} \DataTypeTok{VPrim} \OperatorTok{$} \DataTypeTok{PText}\NormalTok{ (T.replace needle replacement haystack)}

    \CommentTok{{-}{-} Natural/fold: the interesting one}
\NormalTok{    (}\DataTypeTok{NATURAL\_FOLD}\NormalTok{, initVal }\OperatorTok{:|}\NormalTok{ [succ\textquotesingle{}, \_, }\DataTypeTok{VPrim}\NormalTok{ (}\DataTypeTok{PNat}\NormalTok{ n)]) }\OtherTok{{-}\textgreater{}}
\NormalTok{        applyNTimes succ\textquotesingle{} initVal n}

    \CommentTok{{-}{-} List/fold: fold over a list}
\NormalTok{    (}\DataTypeTok{LIST\_FOLD}\NormalTok{, initVal }\OperatorTok{:|}\NormalTok{ [cons, \_, }\DataTypeTok{VList}\NormalTok{ xs, \_]) }\OtherTok{{-}\textgreater{}}
\NormalTok{        V.foldr (\textbackslash{}x acc }\OtherTok{{-}\textgreater{}}\NormalTok{ applyTwice cons x }\OperatorTok{=\textless{}\textless{}}\NormalTok{ acc) (}\FunctionTok{pure}\NormalTok{ initVal) xs}

    \CommentTok{{-}{-} Not saturated yet: accumulate}
\NormalTok{    (builtin\textquotesingle{}, args) }\OtherTok{{-}\textgreater{}} \FunctionTok{pure} \OperatorTok{$} \DataTypeTok{VBuiltin}\NormalTok{ builtin\textquotesingle{} (toList args)}
\end{Highlighting}
\end{Shaded}

\textbf{Natural/fold} applies the successor function N times:

\begin{Shaded}
\begin{Highlighting}[]
\NormalTok{applyNTimes \_ acc }\DecValTok{0} \OtherTok{=} \FunctionTok{pure}\NormalTok{ acc}
\NormalTok{applyNTimes (}\DataTypeTok{VLam}\NormalTok{ procId closedEnv) acc n }\OtherTok{=} \KeywordTok{do}
\NormalTok{    result }\OtherTok{\textless{}{-}}\NormalTok{ executeProcedure procId (acc }\OperatorTok{\textless{}|}\NormalTok{ closedEnv)}
\NormalTok{    applyNTimes (}\DataTypeTok{VLam}\NormalTok{ procId closedEnv) result (n }\OperatorTok{{-}} \DecValTok{1}\NormalTok{)}
\end{Highlighting}
\end{Shaded}

\textbf{List/fold} applies \texttt{cons\ x\ acc} for each element:

\begin{Shaded}
\begin{Highlighting}[]
\NormalTok{applyTwice (}\DataTypeTok{VLam}\NormalTok{ procId1 env1) x acc }\OtherTok{=} \KeywordTok{do}
\NormalTok{    partial }\OtherTok{\textless{}{-}}\NormalTok{ executeProcedure procId1 (x }\OperatorTok{\textless{}|}\NormalTok{ env1)    }\CommentTok{{-}{-} cons x}
    \KeywordTok{case}\NormalTok{ partial }\KeywordTok{of}
        \DataTypeTok{VLam}\NormalTok{ procId2 env2 }\OtherTok{{-}\textgreater{}}\NormalTok{ executeProcedure procId2 (acc }\OperatorTok{\textless{}|}\NormalTok{ env2)  }\CommentTok{{-}{-} (cons x) acc}
\end{Highlighting}
\end{Shaded}

\textbf{Natural/build} and \textbf{List/build} pass runtime-provided functions
to user code:

\begin{Shaded}
\begin{Highlighting}[]
\NormalTok{(}\DataTypeTok{NATURAL\_BUILD}\NormalTok{, }\DataTypeTok{VLam}\NormalTok{ procId closedEnv }\OperatorTok{:|}\NormalTok{ []) }\OtherTok{{-}\textgreater{}} \KeywordTok{do}
    \CommentTok{{-}{-} build f = f Natural Natural/succ 0}
    \CommentTok{{-}{-} We skip the Type argument, pass succ and zero:}
\NormalTok{    partial }\OtherTok{\textless{}{-}}\NormalTok{ executeProcedure procId (}\DataTypeTok{VProvided} \DataTypeTok{NaturalSucc}\NormalTok{ [] }\OperatorTok{\textless{}|}\NormalTok{ closedEnv)}
    \KeywordTok{case}\NormalTok{ partial }\KeywordTok{of}
        \DataTypeTok{VLam}\NormalTok{ procId\textquotesingle{} closedEnv\textquotesingle{} }\OtherTok{{-}\textgreater{}}
\NormalTok{            executeProcedure procId\textquotesingle{} (}\DataTypeTok{VPrim}\NormalTok{ (}\DataTypeTok{PNat} \DecValTok{0}\NormalTok{) }\OperatorTok{\textless{}|}\NormalTok{ closedEnv\textquotesingle{})}
\end{Highlighting}
\end{Shaded}

The \texttt{VProvided\ NaturalSucc} is handed to user code. When user code calls
it (via \texttt{CALL}), \texttt{executeProvidedFn} handles it:

\begin{Shaded}
\begin{Highlighting}[]
\OtherTok{executeProvidedFn ::} \DataTypeTok{ProvidedFn} \OtherTok{{-}\textgreater{}} \DataTypeTok{NonEmpty}\NormalTok{ (}\DataTypeTok{Value} \DataTypeTok{Int}\NormalTok{) }\OtherTok{{-}\textgreater{}} \DataTypeTok{Exec}\NormalTok{ s (}\DataTypeTok{Value} \DataTypeTok{Int}\NormalTok{)}
\NormalTok{executeProvidedFn provFn allArgs }\OtherTok{=} \KeywordTok{case}\NormalTok{ (provFn, allArgs) }\KeywordTok{of}
\NormalTok{    (}\DataTypeTok{NaturalSucc}\NormalTok{, }\DataTypeTok{VPrim}\NormalTok{ (}\DataTypeTok{PNat}\NormalTok{ n) }\OperatorTok{:|}\NormalTok{ []) }\OtherTok{{-}\textgreater{}} \FunctionTok{pure} \OperatorTok{$} \DataTypeTok{VPrim} \OperatorTok{$} \DataTypeTok{PNat}\NormalTok{ (n }\OperatorTok{+} \DecValTok{1}\NormalTok{)}
\NormalTok{    (}\DataTypeTok{ListCons}\NormalTok{, x }\OperatorTok{:|}\NormalTok{ [}\DataTypeTok{VList}\NormalTok{ xs])         }\OtherTok{{-}\textgreater{}} \FunctionTok{pure} \OperatorTok{$} \DataTypeTok{VList}\NormalTok{ (V.cons x xs)}
\NormalTok{    (provFn\textquotesingle{}, args)                      }\OtherTok{{-}\textgreater{}} \FunctionTok{pure} \OperatorTok{$} \DataTypeTok{VProvided}\NormalTok{ provFn\textquotesingle{} (toList args)}
\end{Highlighting}
\end{Shaded}

\subsection{The RunResult protocol}\label{the-runresult-protocol}

\begin{Shaded}
\begin{Highlighting}[]
\KeywordTok{data} \DataTypeTok{RunResult}\NormalTok{ s r a}
    \OtherTok{=} \DataTypeTok{RunFailure} \OperatorTok{!}\NormalTok{(}\DataTypeTok{ExecError}\NormalTok{ s)}
    \OperatorTok{|} \DataTypeTok{RunClosure}\NormalTok{ (r }\OtherTok{{-}\textgreater{}} \DataTypeTok{RunResult}\NormalTok{ s r a)}
    \OperatorTok{|} \DataTypeTok{RunSuccess} \OperatorTok{!}\NormalTok{a}
\end{Highlighting}
\end{Shaded}

When executing a program, if the result is a \texttt{VLam}, it becomes
\texttt{RunClosure} -- a Haskell function waiting for the next argument. This is
the interface for calling compiled Dhall functions from Haskell incrementally:

\begin{Shaded}
\begin{Highlighting}[]
\OtherTok{valueToRunResult ::} \DataTypeTok{ProgramBanks} \OtherTok{{-}\textgreater{}} \DataTypeTok{Value} \DataTypeTok{Int} \OtherTok{{-}\textgreater{}} \DataTypeTok{RunResult}\NormalTok{ s }\DataTypeTok{EValue} \DataTypeTok{EValue}
\NormalTok{valueToRunResult banks (}\DataTypeTok{VLam}\NormalTok{ procId closedEnv) }\OtherTok{=} \DataTypeTok{RunClosure} \OperatorTok{$}\NormalTok{ \textbackslash{}arg }\OtherTok{{-}\textgreater{}}
    \KeywordTok{case}\NormalTok{ evalueToValue arg }\KeywordTok{of}
        \DataTypeTok{Right}\NormalTok{ argVal }\OtherTok{{-}\textgreater{}} \KeywordTok{case}\NormalTok{ runExec (executeProcedure procId (argVal }\OperatorTok{\textless{}|}\NormalTok{ closedEnv)) }\KeywordTok{of}
            \DataTypeTok{Right}\NormalTok{ result }\OtherTok{{-}\textgreater{}}\NormalTok{ valueToRunResult banks result}
            \DataTypeTok{Left}\NormalTok{ err }\OtherTok{{-}\textgreater{}} \DataTypeTok{RunFailure}\NormalTok{ err}
\NormalTok{valueToRunResult \_ (}\DataTypeTok{VPrim}\NormalTok{ p) }\OtherTok{=} \DataTypeTok{RunSuccess} \OperatorTok{$} \DataTypeTok{EPrim}\NormalTok{ p}
\CommentTok{{-}{-} ... (records, lists, unions convert recursively)}
\end{Highlighting}
\end{Shaded}

\texttt{EValue} is the closure-free subset of \texttt{Value} -- the external
interface type:

\begin{Shaded}
\begin{Highlighting}[]
\KeywordTok{data} \DataTypeTok{EValue}
    \OtherTok{=} \DataTypeTok{EPrim} \OperatorTok{!}\DataTypeTok{Prim}
    \OperatorTok{|} \DataTypeTok{ERecord} \OperatorTok{!}\NormalTok{(}\DataTypeTok{Vector} \DataTypeTok{EValue}\NormalTok{)}
    \OperatorTok{|} \DataTypeTok{EList} \OperatorTok{!}\NormalTok{(}\DataTypeTok{Vector} \DataTypeTok{EValue}\NormalTok{)}
    \OperatorTok{|} \DataTypeTok{EUnion} \OperatorTok{!}\DataTypeTok{Word8} \OperatorTok{!}\NormalTok{(}\DataTypeTok{Maybe} \DataTypeTok{EValue}\NormalTok{)}
    \OperatorTok{|} \DataTypeTok{EBuiltin} \OperatorTok{!}\DataTypeTok{DhallBuiltin}\NormalTok{ [}\DataTypeTok{EValue}\NormalTok{]}
    \OperatorTok{|} \DataTypeTok{EProvided} \OperatorTok{!}\DataTypeTok{ProvidedFn}\NormalTok{ [}\DataTypeTok{EValue}\NormalTok{]}
\end{Highlighting}
\end{Shaded}

\section{The C Interpreter
(dhall\_interpreter.c)}\label{the-c-interpreter-dhall_interpreter.c}

A direct translation of the Haskell interpreter into C, operating over the same
bytecode serialized into C structs.

\subsection{Data layout}\label{data-layout}

\begin{Shaded}
\begin{Highlighting}[]
\KeywordTok{typedef} \KeywordTok{enum} \OperatorTok{\{}
\NormalTok{  PRIM\_NAT}\OperatorTok{,}\NormalTok{ PRIM\_INT}\OperatorTok{,}\NormalTok{ PRIM\_DOUBLE}\OperatorTok{,}\NormalTok{ PRIM\_BOOL}\OperatorTok{,}\NormalTok{ PRIM\_TEXT}\OperatorTok{,}
\NormalTok{  PRIM\_BYTES}\OperatorTok{,}\NormalTok{ PRIM\_DATE}\OperatorTok{,}\NormalTok{ PRIM\_TIME}\OperatorTok{,}\NormalTok{ PRIM\_TIMEZONE}\OperatorTok{,}\NormalTok{ PRIM\_ASSERT}\OperatorTok{,}\NormalTok{ PRIM\_TYPE}
\OperatorTok{\}}\NormalTok{ PrimTag}\OperatorTok{;}

\KeywordTok{typedef} \KeywordTok{struct} \OperatorTok{\{}
\NormalTok{  PrimTag tag}\OperatorTok{;}
  \KeywordTok{union} \OperatorTok{\{}
    \DataTypeTok{uint64\_t}\NormalTok{ nat}\OperatorTok{;}
    \DataTypeTok{int64\_t}\NormalTok{ int\_val}\OperatorTok{;}
    \DataTypeTok{double}\NormalTok{ double\_val}\OperatorTok{;}
    \DataTypeTok{bool}\NormalTok{ bool\_val}\OperatorTok{;}
    \DataTypeTok{char} \OperatorTok{*}\NormalTok{text}\OperatorTok{;}
    \KeywordTok{struct} \OperatorTok{\{} \DataTypeTok{uint8\_t} \OperatorTok{*}\NormalTok{data}\OperatorTok{;} \DataTypeTok{size\_t}\NormalTok{ len}\OperatorTok{;} \OperatorTok{\}}\NormalTok{ bytes}\OperatorTok{;}
    \DataTypeTok{int32\_t}\NormalTok{ date}\OperatorTok{;}      \CommentTok{// days since epoch}
    \DataTypeTok{int64\_t}\NormalTok{ time}\OperatorTok{;}      \CommentTok{// nanoseconds since midnight}
    \DataTypeTok{int32\_t}\NormalTok{ timezone}\OperatorTok{;}  \CommentTok{// offset in minutes}
  \OperatorTok{\}}\NormalTok{ data}\OperatorTok{;}
\OperatorTok{\}}\NormalTok{ Prim}\OperatorTok{;}

\KeywordTok{typedef} \KeywordTok{enum} \OperatorTok{\{}
\NormalTok{  VAL\_PRIM}\OperatorTok{,}\NormalTok{ VAL\_RECORD}\OperatorTok{,}\NormalTok{ VAL\_LIST}\OperatorTok{,}\NormalTok{ VAL\_UNION}\OperatorTok{,}
\NormalTok{  VAL\_LAM}\OperatorTok{,}\NormalTok{ VAL\_BUILTIN}\OperatorTok{,}\NormalTok{ VAL\_PROVIDED}
\OperatorTok{\}}\NormalTok{ ValueTag}\OperatorTok{;}

\KeywordTok{struct}\NormalTok{ Value }\OperatorTok{\{}
\NormalTok{  ValueTag tag}\OperatorTok{;}
  \KeywordTok{union} \OperatorTok{\{}
\NormalTok{    Prim prim}\OperatorTok{;}
    \KeywordTok{struct} \OperatorTok{\{}\NormalTok{ Value }\OperatorTok{**}\NormalTok{fields}\OperatorTok{;} \DataTypeTok{size\_t}\NormalTok{ len}\OperatorTok{;} \OperatorTok{\}}\NormalTok{ record}\OperatorTok{;}
    \KeywordTok{struct} \OperatorTok{\{}\NormalTok{ Value }\OperatorTok{**}\NormalTok{elems}\OperatorTok{;} \DataTypeTok{size\_t}\NormalTok{ len}\OperatorTok{;} \OperatorTok{\}}\NormalTok{ list}\OperatorTok{;}
    \KeywordTok{struct} \OperatorTok{\{} \DataTypeTok{uint8\_t}\NormalTok{ tag}\OperatorTok{;}\NormalTok{ Value }\OperatorTok{*}\NormalTok{payload}\OperatorTok{;} \OperatorTok{\}}\NormalTok{ union\_val}\OperatorTok{;}
    \KeywordTok{struct} \OperatorTok{\{} \DataTypeTok{int}\NormalTok{ proc\_id}\OperatorTok{;}\NormalTok{ Value }\OperatorTok{**}\NormalTok{env}\OperatorTok{;} \DataTypeTok{size\_t}\NormalTok{ env\_len}\OperatorTok{;} \OperatorTok{\}}\NormalTok{ lam}\OperatorTok{;}
    \KeywordTok{struct} \OperatorTok{\{}\NormalTok{ DhallBuiltin builtin}\OperatorTok{;}\NormalTok{ Value }\OperatorTok{**}\NormalTok{args}\OperatorTok{;} \DataTypeTok{size\_t}\NormalTok{ args\_len}\OperatorTok{;} \OperatorTok{\}}\NormalTok{ builtin}\OperatorTok{;}
    \KeywordTok{struct} \OperatorTok{\{}\NormalTok{ ProvidedFn fn}\OperatorTok{;}\NormalTok{ Value }\OperatorTok{**}\NormalTok{args}\OperatorTok{;} \DataTypeTok{size\_t}\NormalTok{ args\_len}\OperatorTok{;} \OperatorTok{\}}\NormalTok{ provided}\OperatorTok{;}
  \OperatorTok{\}}\NormalTok{ data}\OperatorTok{;}
\OperatorTok{\};}
\end{Highlighting}
\end{Shaded}

\subsection{Program structure}\label{program-structure}

\begin{Shaded}
\begin{Highlighting}[]
\KeywordTok{typedef} \KeywordTok{struct} \OperatorTok{\{}
  \DataTypeTok{uint8\_t}\NormalTok{ tag}\OperatorTok{;}  \CommentTok{// InstrTag}
  \KeywordTok{union} \OperatorTok{\{}
    \DataTypeTok{uint8\_t}\NormalTok{ load\_local\_idx}\OperatorTok{;}
    \DataTypeTok{uint16\_t}\NormalTok{ make\_list\_n}\OperatorTok{;}
    \DataTypeTok{uint16\_t}\NormalTok{ make\_record\_n}\OperatorTok{;}
    \DataTypeTok{uint16\_t}\NormalTok{ load\_field\_idx}\OperatorTok{;}
    \DataTypeTok{uint8\_t}\NormalTok{ make\_union\_tag}\OperatorTok{;}
    \DataTypeTok{uint8\_t}\NormalTok{ make\_nullary\_union\_tag}\OperatorTok{;}
    \KeywordTok{struct} \OperatorTok{\{}\NormalTok{ UnionArity }\OperatorTok{*}\NormalTok{arities}\OperatorTok{;} \DataTypeTok{size\_t}\NormalTok{ arities\_len}\OperatorTok{;} \OperatorTok{\}}\NormalTok{ merge\_union}\OperatorTok{;}
    \KeywordTok{struct} \OperatorTok{\{} \DataTypeTok{int}\NormalTok{ true\_proc}\OperatorTok{;} \DataTypeTok{int}\NormalTok{ false\_proc}\OperatorTok{;} \OperatorTok{\}}\NormalTok{ merge\_if}\OperatorTok{;}
\NormalTok{    Prim push\_prim}\OperatorTok{;}
\NormalTok{    DhallBuiltin push\_builtin}\OperatorTok{;}
\NormalTok{    DhallBuiltin binop}\OperatorTok{;}
    \DataTypeTok{int}\NormalTok{ closure\_proc}\OperatorTok{;}
    \KeywordTok{struct} \OperatorTok{\{} \DataTypeTok{char} \OperatorTok{**}\NormalTok{names}\OperatorTok{;} \DataTypeTok{size\_t}\NormalTok{ names\_len}\OperatorTok{;} \OperatorTok{\}}\NormalTok{ show\_constructor}\OperatorTok{;}
    \KeywordTok{struct} \OperatorTok{\{} \DataTypeTok{char} \OperatorTok{**}\NormalTok{names}\OperatorTok{;} \DataTypeTok{size\_t}\NormalTok{ names\_len}\OperatorTok{;} \OperatorTok{\}}\NormalTok{ tomap}\OperatorTok{;}
    \KeywordTok{struct} \OperatorTok{\{}\NormalTok{ PathComponent }\OperatorTok{*}\NormalTok{path}\OperatorTok{;} \DataTypeTok{size\_t}\NormalTok{ path\_len}\OperatorTok{;} \OperatorTok{\}}\NormalTok{ with}\OperatorTok{;}
    \KeywordTok{struct} \OperatorTok{\{} \DataTypeTok{uint16\_t} \OperatorTok{*}\NormalTok{indices}\OperatorTok{;} \DataTypeTok{size\_t}\NormalTok{ indices\_len}\OperatorTok{;} \OperatorTok{\}}\NormalTok{ project}\OperatorTok{;}
    \KeywordTok{struct} \OperatorTok{\{}\NormalTok{ RecordMergeOpSeq }\OperatorTok{*}\NormalTok{ops}\OperatorTok{;} \DataTypeTok{size\_t}\NormalTok{ ops\_len}\OperatorTok{;} \OperatorTok{\}}\NormalTok{ combine}\OperatorTok{;}
    \KeywordTok{struct} \OperatorTok{\{}\NormalTok{ RecordMergeOp }\OperatorTok{*}\NormalTok{ops}\OperatorTok{;} \DataTypeTok{size\_t}\NormalTok{ ops\_len}\OperatorTok{;} \OperatorTok{\}}\NormalTok{ prefer}\OperatorTok{;}
  \OperatorTok{\}}\NormalTok{ data}\OperatorTok{;}
\OperatorTok{\}}\NormalTok{ Instr}\OperatorTok{;}

\KeywordTok{typedef} \KeywordTok{struct} \OperatorTok{\{}\NormalTok{ Instr }\OperatorTok{*}\NormalTok{instrs}\OperatorTok{;} \DataTypeTok{size\_t}\NormalTok{ len}\OperatorTok{;} \OperatorTok{\}}\NormalTok{ Procedure}\OperatorTok{;}
\KeywordTok{typedef} \KeywordTok{struct} \OperatorTok{\{}\NormalTok{ Procedure }\OperatorTok{*}\NormalTok{procs}\OperatorTok{;} \DataTypeTok{size\_t}\NormalTok{ num\_procs}\OperatorTok{;} \OperatorTok{\}}\NormalTok{ Program}\OperatorTok{;}
\end{Highlighting}
\end{Shaded}

\subsection{Execution context}\label{execution-context}

\begin{Shaded}
\begin{Highlighting}[]
\PreprocessorTok{\#define MAX\_STACK\_SIZE }\DecValTok{1024}

\KeywordTok{struct}\NormalTok{ ExecutionContext }\OperatorTok{\{}
\NormalTok{  Value }\OperatorTok{*}\NormalTok{stack}\OperatorTok{[}\NormalTok{MAX\_STACK\_SIZE}\OperatorTok{];}
  \DataTypeTok{size\_t}\NormalTok{ stack\_len}\OperatorTok{;}
\NormalTok{  Value }\OperatorTok{**}\NormalTok{env}\OperatorTok{;}
  \DataTypeTok{size\_t}\NormalTok{ env\_len}\OperatorTok{;}
\NormalTok{  Program }\OperatorTok{*}\NormalTok{program}\OperatorTok{;}
\OperatorTok{\};}
\end{Highlighting}
\end{Shaded}

Fixed-size stack array. The interpreter pushes and pops pointers to heap-
allocated Values.

\subsection{Core execution}\label{core-execution}

\begin{Shaded}
\begin{Highlighting}[]
\NormalTok{Value }\OperatorTok{*}\NormalTok{execute\_procedure}\OperatorTok{(}\NormalTok{ExecutionContext }\OperatorTok{*}\NormalTok{ctx}\OperatorTok{,} \DataTypeTok{int}\NormalTok{ proc\_id}\OperatorTok{,}
\NormalTok{                         Value }\OperatorTok{**}\NormalTok{env}\OperatorTok{,} \DataTypeTok{size\_t}\NormalTok{ env\_len}\OperatorTok{,}\NormalTok{ ErrorCode }\OperatorTok{*}\NormalTok{err}\OperatorTok{)} \OperatorTok{\{}
    \CommentTok{// Save/restore env around the call}
\NormalTok{    Value }\OperatorTok{**}\NormalTok{old\_env }\OperatorTok{=}\NormalTok{ ctx}\OperatorTok{{-}\textgreater{}}\NormalTok{env}\OperatorTok{;}
    \DataTypeTok{size\_t}\NormalTok{ old\_env\_len }\OperatorTok{=}\NormalTok{ ctx}\OperatorTok{{-}\textgreater{}}\NormalTok{env\_len}\OperatorTok{;}
\NormalTok{    ctx}\OperatorTok{{-}\textgreater{}}\NormalTok{env }\OperatorTok{=}\NormalTok{ env}\OperatorTok{;}
\NormalTok{    ctx}\OperatorTok{{-}\textgreater{}}\NormalTok{env\_len }\OperatorTok{=}\NormalTok{ env\_len}\OperatorTok{;}

\NormalTok{    Procedure }\OperatorTok{*}\NormalTok{proc }\OperatorTok{=} \OperatorTok{\&}\NormalTok{ctx}\OperatorTok{{-}\textgreater{}}\NormalTok{program}\OperatorTok{{-}\textgreater{}}\NormalTok{procs}\OperatorTok{[}\NormalTok{proc\_id}\OperatorTok{];}
    \ControlFlowTok{for} \OperatorTok{(}\DataTypeTok{size\_t}\NormalTok{ i }\OperatorTok{=} \DecValTok{0}\OperatorTok{;}\NormalTok{ i }\OperatorTok{\textless{}}\NormalTok{ proc}\OperatorTok{{-}\textgreater{}}\NormalTok{len }\OperatorTok{\&\&} \OperatorTok{*}\NormalTok{err }\OperatorTok{==}\NormalTok{ ERR\_NONE}\OperatorTok{;}\NormalTok{ i}\OperatorTok{++)} \OperatorTok{\{}
\NormalTok{        execute\_instr}\OperatorTok{(}\NormalTok{ctx}\OperatorTok{,} \OperatorTok{\&}\NormalTok{proc}\OperatorTok{{-}\textgreater{}}\NormalTok{instrs}\OperatorTok{[}\NormalTok{i}\OperatorTok{],}\NormalTok{ err}\OperatorTok{);}
    \OperatorTok{\}}

\NormalTok{    ctx}\OperatorTok{{-}\textgreater{}}\NormalTok{env }\OperatorTok{=}\NormalTok{ old\_env}\OperatorTok{;}
\NormalTok{    ctx}\OperatorTok{{-}\textgreater{}}\NormalTok{env\_len }\OperatorTok{=}\NormalTok{ old\_env\_len}\OperatorTok{;}
    \ControlFlowTok{return}\NormalTok{ pop}\OperatorTok{(}\NormalTok{ctx}\OperatorTok{,}\NormalTok{ err}\OperatorTok{);}
\OperatorTok{\}}
\end{Highlighting}
\end{Shaded}

Each instruction case in \texttt{execute\_instr} mirrors the Haskell interpreter
exactly -- CLOSURE allocates a Value with tag VAL\_LAM and memcpy's the env,
CALL checks the function tag and either recurses into execute\_procedure (for
closures) or accumulates args (for builtins), etc.

\subsection{CALL in C}\label{call-in-c}

\begin{Shaded}
\begin{Highlighting}[]
\ControlFlowTok{case}\NormalTok{ INSTR\_CALL}\OperatorTok{:} \OperatorTok{\{}
\NormalTok{    Value }\OperatorTok{*}\NormalTok{arg }\OperatorTok{=}\NormalTok{ pop}\OperatorTok{(}\NormalTok{ctx}\OperatorTok{,}\NormalTok{ err}\OperatorTok{);}
\NormalTok{    Value }\OperatorTok{*}\NormalTok{func }\OperatorTok{=}\NormalTok{ pop}\OperatorTok{(}\NormalTok{ctx}\OperatorTok{,}\NormalTok{ err}\OperatorTok{);}

    \ControlFlowTok{if} \OperatorTok{(}\NormalTok{func}\OperatorTok{{-}\textgreater{}}\NormalTok{tag }\OperatorTok{==}\NormalTok{ VAL\_LAM}\OperatorTok{)} \OperatorTok{\{}
        \DataTypeTok{size\_t}\NormalTok{ new\_env\_len }\OperatorTok{=}\NormalTok{ func}\OperatorTok{{-}\textgreater{}}\NormalTok{data}\OperatorTok{.}\NormalTok{lam}\OperatorTok{.}\NormalTok{env\_len }\OperatorTok{+} \DecValTok{1}\OperatorTok{;}
\NormalTok{        Value }\OperatorTok{**}\NormalTok{new\_env }\OperatorTok{=}\NormalTok{ malloc}\OperatorTok{(}\KeywordTok{sizeof}\OperatorTok{(}\NormalTok{Value }\OperatorTok{*)} \OperatorTok{*}\NormalTok{ new\_env\_len}\OperatorTok{);}
\NormalTok{        new\_env}\OperatorTok{[}\DecValTok{0}\OperatorTok{]} \OperatorTok{=}\NormalTok{ arg}\OperatorTok{;}
\NormalTok{        memcpy}\OperatorTok{(}\NormalTok{new\_env }\OperatorTok{+} \DecValTok{1}\OperatorTok{,}\NormalTok{ func}\OperatorTok{{-}\textgreater{}}\NormalTok{data}\OperatorTok{.}\NormalTok{lam}\OperatorTok{.}\NormalTok{env}\OperatorTok{,}
\NormalTok{               func}\OperatorTok{{-}\textgreater{}}\NormalTok{data}\OperatorTok{.}\NormalTok{lam}\OperatorTok{.}\NormalTok{env\_len }\OperatorTok{*} \KeywordTok{sizeof}\OperatorTok{(}\NormalTok{Value }\OperatorTok{*));}
\NormalTok{        Value }\OperatorTok{*}\NormalTok{result }\OperatorTok{=}\NormalTok{ execute\_procedure}\OperatorTok{(}\NormalTok{ctx}\OperatorTok{,}\NormalTok{ func}\OperatorTok{{-}\textgreater{}}\NormalTok{data}\OperatorTok{.}\NormalTok{lam}\OperatorTok{.}\NormalTok{proc\_id}\OperatorTok{,}
\NormalTok{                                          new\_env}\OperatorTok{,}\NormalTok{ new\_env\_len}\OperatorTok{,}\NormalTok{ err}\OperatorTok{);}
\NormalTok{        push}\OperatorTok{(}\NormalTok{ctx}\OperatorTok{,}\NormalTok{ result}\OperatorTok{,}\NormalTok{ err}\OperatorTok{);}
    \OperatorTok{\}} \ControlFlowTok{else} \ControlFlowTok{if} \OperatorTok{(}\NormalTok{func}\OperatorTok{{-}\textgreater{}}\NormalTok{tag }\OperatorTok{==}\NormalTok{ VAL\_BUILTIN}\OperatorTok{)} \OperatorTok{\{}
        \CommentTok{// Accumulate arg; if saturated, fire builtin}
        \OperatorTok{...}
    \OperatorTok{\}}
    \ControlFlowTok{break}\OperatorTok{;}
\OperatorTok{\}}
\end{Highlighting}
\end{Shaded}

\subsection{The FFI boundary}\label{the-ffi-boundary}

Haskell serializes \texttt{ProgramBanks} to a binary format, passes it across
FFI as bytes, C deserializes into a \texttt{Program} struct, executes, and
returns an \texttt{ExecutionResult}:

\begin{Shaded}
\begin{Highlighting}[]
\KeywordTok{typedef} \KeywordTok{enum} \OperatorTok{\{}\NormalTok{ RESULT\_SUCCESS}\OperatorTok{,}\NormalTok{ RESULT\_CLOSURE}\OperatorTok{,}\NormalTok{ RESULT\_FAILURE }\OperatorTok{\}}\NormalTok{ ResultTag}\OperatorTok{;}

\KeywordTok{typedef} \KeywordTok{struct} \OperatorTok{\{}
\NormalTok{  ResultTag tag}\OperatorTok{;}
  \KeywordTok{union} \OperatorTok{\{}
\NormalTok{    Value }\OperatorTok{*}\NormalTok{success}\OperatorTok{;}
    \KeywordTok{struct} \OperatorTok{\{} \DataTypeTok{int}\NormalTok{ proc\_id}\OperatorTok{;}\NormalTok{ Value }\OperatorTok{**}\NormalTok{captured\_env}\OperatorTok{;} \DataTypeTok{size\_t}\NormalTok{ env\_len}\OperatorTok{;} \OperatorTok{\}}\NormalTok{ closure}\OperatorTok{;}
    \KeywordTok{struct} \OperatorTok{\{}\NormalTok{ ErrorCode code}\OperatorTok{;} \DataTypeTok{char} \OperatorTok{*}\NormalTok{message}\OperatorTok{;} \OperatorTok{\}}\NormalTok{ error}\OperatorTok{;}
  \OperatorTok{\}}\NormalTok{ data}\OperatorTok{;}
\OperatorTok{\}}\NormalTok{ ExecutionResult}\OperatorTok{;}
\end{Highlighting}
\end{Shaded}

If the result is \texttt{RESULT\_CLOSURE}, the Haskell side holds onto the
proc\_id and captured\_env, and can feed more arguments in subsequent FFI calls.
This is the progressive application protocol -- the same as \texttt{RunClosure}
on the Haskell side.

\section{The Native Backend (Compiled C via
GCC)}\label{the-native-backend-compiled-c-via-gcc}

The native backend translates bytecode directly to C source, compiles it with
GCC -O2, and dlopens the resulting \texttt{.so}. Each procedure bank becomes a
static C function; inter-bank calls go through a \texttt{dispatch} switch.

Given this Dhall expression:

\begin{Shaded}
\begin{Highlighting}[]
\OperatorTok{\textbackslash{}}\NormalTok{(xs}\CommentTok{ }\NormalTok{:}\CommentTok{ }\NormalTok{List}\CommentTok{ }\NormalTok{Natural)}\CommentTok{ }\OperatorTok{{-}\textgreater{}}
\CommentTok{  }\NormalTok{List/fold}\CommentTok{ }\NormalTok{Natural}\CommentTok{ }\NormalTok{xs}\CommentTok{ }\NormalTok{Natural}
\CommentTok{    }\NormalTok{(}\OperatorTok{\textbackslash{}}\NormalTok{(x}\CommentTok{ }\NormalTok{:}\CommentTok{ }\NormalTok{Natural)}\CommentTok{ }\OperatorTok{{-}\textgreater{}}\CommentTok{ }\OperatorTok{\textbackslash{}}\NormalTok{(acc}\CommentTok{ }\NormalTok{:}\CommentTok{ }\NormalTok{Natural)}\CommentTok{ }\OperatorTok{{-}\textgreater{}}\CommentTok{ }\NormalTok{x}\CommentTok{ }\NormalTok{+}\CommentTok{ }\NormalTok{acc)}\CommentTok{ }\DecValTok{0}
\end{Highlighting}
\end{Shaded}

The native backend emits (literal output of
\texttt{dhall-jit-print\ -\/-csource}):

\begin{Shaded}
\begin{Highlighting}[]
\PreprocessorTok{\#include }\ImportTok{\textless{}dhall\_native\_rt.h\textgreater{}}

\DataTypeTok{static}\NormalTok{ Value }\OperatorTok{*}\NormalTok{proc\_0}\OperatorTok{(}\NormalTok{Arena }\OperatorTok{*}\NormalTok{a}\OperatorTok{,}\NormalTok{ Value }\OperatorTok{**}\NormalTok{env}\OperatorTok{,} \DataTypeTok{size\_t}\NormalTok{ env\_len}\OperatorTok{);}
\DataTypeTok{static}\NormalTok{ Value }\OperatorTok{*}\NormalTok{proc\_1}\OperatorTok{(}\NormalTok{Arena }\OperatorTok{*}\NormalTok{a}\OperatorTok{,}\NormalTok{ Value }\OperatorTok{**}\NormalTok{env}\OperatorTok{,} \DataTypeTok{size\_t}\NormalTok{ env\_len}\OperatorTok{);}
\DataTypeTok{static}\NormalTok{ Value }\OperatorTok{*}\NormalTok{proc\_2}\OperatorTok{(}\NormalTok{Arena }\OperatorTok{*}\NormalTok{a}\OperatorTok{,}\NormalTok{ Value }\OperatorTok{**}\NormalTok{env}\OperatorTok{,} \DataTypeTok{size\_t}\NormalTok{ env\_len}\OperatorTok{);}
\DataTypeTok{static}\NormalTok{ Value }\OperatorTok{*}\NormalTok{proc\_3}\OperatorTok{(}\NormalTok{Arena }\OperatorTok{*}\NormalTok{a}\OperatorTok{,}\NormalTok{ Value }\OperatorTok{**}\NormalTok{env}\OperatorTok{,} \DataTypeTok{size\_t}\NormalTok{ env\_len}\OperatorTok{);}

\DataTypeTok{static}\NormalTok{ Value }\OperatorTok{*}\NormalTok{proc\_0}\OperatorTok{(}\NormalTok{Arena }\OperatorTok{*}\NormalTok{a}\OperatorTok{,}\NormalTok{ Value }\OperatorTok{**}\NormalTok{env}\OperatorTok{,} \DataTypeTok{size\_t}\NormalTok{ env\_len}\OperatorTok{)} \OperatorTok{\{}
\NormalTok{    Value }\OperatorTok{*}\NormalTok{stack}\OperatorTok{[}\DecValTok{1024}\OperatorTok{];}
    \DataTypeTok{int}\NormalTok{ sp }\OperatorTok{=} \DecValTok{0}\OperatorTok{;}
    \OperatorTok{\{}
\NormalTok{        Value }\OperatorTok{*}\NormalTok{lam }\OperatorTok{=}\NormalTok{ value\_alloc}\OperatorTok{(}\NormalTok{a}\OperatorTok{);}
\NormalTok{        lam}\OperatorTok{{-}\textgreater{}}\NormalTok{tag }\OperatorTok{=}\NormalTok{ VAL\_LAM}\OperatorTok{;}
\NormalTok{        lam}\OperatorTok{{-}\textgreater{}}\NormalTok{data}\OperatorTok{.}\NormalTok{lam}\OperatorTok{.}\NormalTok{proc\_id }\OperatorTok{=} \DecValTok{1}\OperatorTok{;}
\NormalTok{        lam}\OperatorTok{{-}\textgreater{}}\NormalTok{data}\OperatorTok{.}\NormalTok{lam}\OperatorTok{.}\NormalTok{env\_len }\OperatorTok{=}\NormalTok{ env\_len}\OperatorTok{;}
        \ControlFlowTok{if} \OperatorTok{(}\NormalTok{env\_len }\OperatorTok{\textgreater{}} \DecValTok{0}\OperatorTok{)} \OperatorTok{\{}
\NormalTok{            lam}\OperatorTok{{-}\textgreater{}}\NormalTok{data}\OperatorTok{.}\NormalTok{lam}\OperatorTok{.}\NormalTok{env }\OperatorTok{=}\NormalTok{ arena\_alloc}\OperatorTok{(}\NormalTok{a}\OperatorTok{,}\NormalTok{ env\_len }\OperatorTok{*} \KeywordTok{sizeof}\OperatorTok{(}\NormalTok{Value}\OperatorTok{*));}
\NormalTok{            memcpy}\OperatorTok{(}\NormalTok{lam}\OperatorTok{{-}\textgreater{}}\NormalTok{data}\OperatorTok{.}\NormalTok{lam}\OperatorTok{.}\NormalTok{env}\OperatorTok{,}\NormalTok{ env}\OperatorTok{,}\NormalTok{ env\_len }\OperatorTok{*} \KeywordTok{sizeof}\OperatorTok{(}\NormalTok{Value}\OperatorTok{*));}
        \OperatorTok{\}} \ControlFlowTok{else} \OperatorTok{\{}
\NormalTok{            lam}\OperatorTok{{-}\textgreater{}}\NormalTok{data}\OperatorTok{.}\NormalTok{lam}\OperatorTok{.}\NormalTok{env }\OperatorTok{=}\NormalTok{ NULL}\OperatorTok{;}
        \OperatorTok{\}}
\NormalTok{        stack}\OperatorTok{[}\NormalTok{sp}\OperatorTok{++]} \OperatorTok{=}\NormalTok{ lam}\OperatorTok{;}
    \OperatorTok{\}}
    \ControlFlowTok{return}\NormalTok{ stack}\OperatorTok{[}\NormalTok{sp }\OperatorTok{{-}} \DecValTok{1}\OperatorTok{];}
\OperatorTok{\}}

\DataTypeTok{static}\NormalTok{ Value }\OperatorTok{*}\NormalTok{proc\_1}\OperatorTok{(}\NormalTok{Arena }\OperatorTok{*}\NormalTok{a}\OperatorTok{,}\NormalTok{ Value }\OperatorTok{**}\NormalTok{env}\OperatorTok{,} \DataTypeTok{size\_t}\NormalTok{ env\_len}\OperatorTok{)} \OperatorTok{\{}
\NormalTok{    Value }\OperatorTok{*}\NormalTok{stack}\OperatorTok{[}\DecValTok{1024}\OperatorTok{];}
    \DataTypeTok{int}\NormalTok{ sp }\OperatorTok{=} \DecValTok{0}\OperatorTok{;}
\NormalTok{    stack}\OperatorTok{[}\NormalTok{sp}\OperatorTok{++]} \OperatorTok{=}\NormalTok{ env}\OperatorTok{[}\DecValTok{0}\OperatorTok{];}
    \OperatorTok{\{}
\NormalTok{        Value }\OperatorTok{*}\NormalTok{lam }\OperatorTok{=}\NormalTok{ value\_alloc}\OperatorTok{(}\NormalTok{a}\OperatorTok{);}
\NormalTok{        lam}\OperatorTok{{-}\textgreater{}}\NormalTok{tag }\OperatorTok{=}\NormalTok{ VAL\_LAM}\OperatorTok{;}
\NormalTok{        lam}\OperatorTok{{-}\textgreater{}}\NormalTok{data}\OperatorTok{.}\NormalTok{lam}\OperatorTok{.}\NormalTok{proc\_id }\OperatorTok{=} \DecValTok{2}\OperatorTok{;}
\NormalTok{        lam}\OperatorTok{{-}\textgreater{}}\NormalTok{data}\OperatorTok{.}\NormalTok{lam}\OperatorTok{.}\NormalTok{env\_len }\OperatorTok{=}\NormalTok{ env\_len}\OperatorTok{;}
        \ControlFlowTok{if} \OperatorTok{(}\NormalTok{env\_len }\OperatorTok{\textgreater{}} \DecValTok{0}\OperatorTok{)} \OperatorTok{\{}
\NormalTok{            lam}\OperatorTok{{-}\textgreater{}}\NormalTok{data}\OperatorTok{.}\NormalTok{lam}\OperatorTok{.}\NormalTok{env }\OperatorTok{=}\NormalTok{ arena\_alloc}\OperatorTok{(}\NormalTok{a}\OperatorTok{,}\NormalTok{ env\_len }\OperatorTok{*} \KeywordTok{sizeof}\OperatorTok{(}\NormalTok{Value}\OperatorTok{*));}
\NormalTok{            memcpy}\OperatorTok{(}\NormalTok{lam}\OperatorTok{{-}\textgreater{}}\NormalTok{data}\OperatorTok{.}\NormalTok{lam}\OperatorTok{.}\NormalTok{env}\OperatorTok{,}\NormalTok{ env}\OperatorTok{,}\NormalTok{ env\_len }\OperatorTok{*} \KeywordTok{sizeof}\OperatorTok{(}\NormalTok{Value}\OperatorTok{*));}
        \OperatorTok{\}} \ControlFlowTok{else} \OperatorTok{\{}
\NormalTok{            lam}\OperatorTok{{-}\textgreater{}}\NormalTok{data}\OperatorTok{.}\NormalTok{lam}\OperatorTok{.}\NormalTok{env }\OperatorTok{=}\NormalTok{ NULL}\OperatorTok{;}
        \OperatorTok{\}}
\NormalTok{        stack}\OperatorTok{[}\NormalTok{sp}\OperatorTok{++]} \OperatorTok{=}\NormalTok{ lam}\OperatorTok{;}
    \OperatorTok{\}}
\NormalTok{    stack}\OperatorTok{[}\NormalTok{sp}\OperatorTok{++]} \OperatorTok{=}\NormalTok{ make\_prim\_nat\_ui}\OperatorTok{(}\NormalTok{a}\OperatorTok{,} \DecValTok{0}\BuiltInTok{UL}\OperatorTok{);}
    \OperatorTok{\{}\NormalTok{ Value }\OperatorTok{*}\NormalTok{init\_v}\OperatorTok{=}\NormalTok{stack}\OperatorTok{[{-}{-}}\NormalTok{sp}\OperatorTok{];}\NormalTok{ Value }\OperatorTok{*}\NormalTok{cons\_v}\OperatorTok{=}\NormalTok{stack}\OperatorTok{[{-}{-}}\NormalTok{sp}\OperatorTok{];}\NormalTok{ Value }\OperatorTok{*}\NormalTok{list\_v}\OperatorTok{=}\NormalTok{stack}\OperatorTok{[{-}{-}}\NormalTok{sp}\OperatorTok{];}
      \DataTypeTok{size\_t}\NormalTok{ len}\OperatorTok{=}\NormalTok{list\_v}\OperatorTok{{-}\textgreater{}}\NormalTok{data}\OperatorTok{.}\NormalTok{list}\OperatorTok{.}\NormalTok{len}\OperatorTok{;}
\NormalTok{      Value }\OperatorTok{*}\NormalTok{acc}\OperatorTok{=}\NormalTok{init\_v}\OperatorTok{;}
      \ControlFlowTok{for}\OperatorTok{(}\DataTypeTok{size\_t}\NormalTok{ i}\OperatorTok{=}\NormalTok{len}\OperatorTok{;}\NormalTok{i}\OperatorTok{\textgreater{}}\DecValTok{0}\OperatorTok{;}\NormalTok{i}\OperatorTok{{-}{-})\{}
\NormalTok{        Value }\OperatorTok{*}\NormalTok{elem}\OperatorTok{=}\NormalTok{list\_v}\OperatorTok{{-}\textgreater{}}\NormalTok{data}\OperatorTok{.}\NormalTok{list}\OperatorTok{.}\NormalTok{elems}\OperatorTok{[}\NormalTok{i}\OperatorTok{{-}}\DecValTok{1}\OperatorTok{];}
\NormalTok{        Value }\OperatorTok{*}\NormalTok{partial}\OperatorTok{=}\NormalTok{call\_value}\OperatorTok{(}\NormalTok{a}\OperatorTok{,}\NormalTok{ cons\_v}\OperatorTok{,}\NormalTok{ elem}\OperatorTok{);}
\NormalTok{        acc}\OperatorTok{=}\NormalTok{call\_value}\OperatorTok{(}\NormalTok{a}\OperatorTok{,}\NormalTok{ partial}\OperatorTok{,}\NormalTok{ acc}\OperatorTok{);}
      \OperatorTok{\}}
\NormalTok{      stack}\OperatorTok{[}\NormalTok{sp}\OperatorTok{++]=}\NormalTok{acc}\OperatorTok{;} \OperatorTok{\}}
    \ControlFlowTok{return}\NormalTok{ stack}\OperatorTok{[}\NormalTok{sp }\OperatorTok{{-}} \DecValTok{1}\OperatorTok{];}
\OperatorTok{\}}

\DataTypeTok{static}\NormalTok{ Value }\OperatorTok{*}\NormalTok{proc\_2}\OperatorTok{(}\NormalTok{Arena }\OperatorTok{*}\NormalTok{a}\OperatorTok{,}\NormalTok{ Value }\OperatorTok{**}\NormalTok{env}\OperatorTok{,} \DataTypeTok{size\_t}\NormalTok{ env\_len}\OperatorTok{)} \OperatorTok{\{}
\NormalTok{    Value }\OperatorTok{*}\NormalTok{stack}\OperatorTok{[}\DecValTok{1024}\OperatorTok{];}
    \DataTypeTok{int}\NormalTok{ sp }\OperatorTok{=} \DecValTok{0}\OperatorTok{;}
    \OperatorTok{\{}
\NormalTok{        Value }\OperatorTok{*}\NormalTok{lam }\OperatorTok{=}\NormalTok{ value\_alloc}\OperatorTok{(}\NormalTok{a}\OperatorTok{);}
\NormalTok{        lam}\OperatorTok{{-}\textgreater{}}\NormalTok{tag }\OperatorTok{=}\NormalTok{ VAL\_LAM}\OperatorTok{;}
\NormalTok{        lam}\OperatorTok{{-}\textgreater{}}\NormalTok{data}\OperatorTok{.}\NormalTok{lam}\OperatorTok{.}\NormalTok{proc\_id }\OperatorTok{=} \DecValTok{3}\OperatorTok{;}
\NormalTok{        lam}\OperatorTok{{-}\textgreater{}}\NormalTok{data}\OperatorTok{.}\NormalTok{lam}\OperatorTok{.}\NormalTok{env\_len }\OperatorTok{=}\NormalTok{ env\_len}\OperatorTok{;}
        \ControlFlowTok{if} \OperatorTok{(}\NormalTok{env\_len }\OperatorTok{\textgreater{}} \DecValTok{0}\OperatorTok{)} \OperatorTok{\{}
\NormalTok{            lam}\OperatorTok{{-}\textgreater{}}\NormalTok{data}\OperatorTok{.}\NormalTok{lam}\OperatorTok{.}\NormalTok{env }\OperatorTok{=}\NormalTok{ arena\_alloc}\OperatorTok{(}\NormalTok{a}\OperatorTok{,}\NormalTok{ env\_len }\OperatorTok{*} \KeywordTok{sizeof}\OperatorTok{(}\NormalTok{Value}\OperatorTok{*));}
\NormalTok{            memcpy}\OperatorTok{(}\NormalTok{lam}\OperatorTok{{-}\textgreater{}}\NormalTok{data}\OperatorTok{.}\NormalTok{lam}\OperatorTok{.}\NormalTok{env}\OperatorTok{,}\NormalTok{ env}\OperatorTok{,}\NormalTok{ env\_len }\OperatorTok{*} \KeywordTok{sizeof}\OperatorTok{(}\NormalTok{Value}\OperatorTok{*));}
        \OperatorTok{\}} \ControlFlowTok{else} \OperatorTok{\{}
\NormalTok{            lam}\OperatorTok{{-}\textgreater{}}\NormalTok{data}\OperatorTok{.}\NormalTok{lam}\OperatorTok{.}\NormalTok{env }\OperatorTok{=}\NormalTok{ NULL}\OperatorTok{;}
        \OperatorTok{\}}
\NormalTok{        stack}\OperatorTok{[}\NormalTok{sp}\OperatorTok{++]} \OperatorTok{=}\NormalTok{ lam}\OperatorTok{;}
    \OperatorTok{\}}
    \ControlFlowTok{return}\NormalTok{ stack}\OperatorTok{[}\NormalTok{sp }\OperatorTok{{-}} \DecValTok{1}\OperatorTok{];}
\OperatorTok{\}}

\DataTypeTok{static}\NormalTok{ Value }\OperatorTok{*}\NormalTok{proc\_3}\OperatorTok{(}\NormalTok{Arena }\OperatorTok{*}\NormalTok{a}\OperatorTok{,}\NormalTok{ Value }\OperatorTok{**}\NormalTok{env}\OperatorTok{,} \DataTypeTok{size\_t}\NormalTok{ env\_len}\OperatorTok{)} \OperatorTok{\{}
\NormalTok{    Value }\OperatorTok{*}\NormalTok{stack}\OperatorTok{[}\DecValTok{1024}\OperatorTok{];}
    \DataTypeTok{int}\NormalTok{ sp }\OperatorTok{=} \DecValTok{0}\OperatorTok{;}
\NormalTok{    stack}\OperatorTok{[}\NormalTok{sp}\OperatorTok{++]} \OperatorTok{=}\NormalTok{ env}\OperatorTok{[}\DecValTok{1}\OperatorTok{];}
\NormalTok{    stack}\OperatorTok{[}\NormalTok{sp}\OperatorTok{++]} \OperatorTok{=}\NormalTok{ env}\OperatorTok{[}\DecValTok{0}\OperatorTok{];}
    \OperatorTok{\{}\NormalTok{ Value }\OperatorTok{*}\NormalTok{b2}\OperatorTok{=}\NormalTok{stack}\OperatorTok{[{-}{-}}\NormalTok{sp}\OperatorTok{];}\NormalTok{ Value }\OperatorTok{*}\NormalTok{a2}\OperatorTok{=}\NormalTok{stack}\OperatorTok{[{-}{-}}\NormalTok{sp}\OperatorTok{];}
\NormalTok{      stack}\OperatorTok{[}\NormalTok{sp}\OperatorTok{++]=}\NormalTok{op\_add\_nat}\OperatorTok{(}\NormalTok{a}\OperatorTok{,} \OperatorTok{\&}\NormalTok{a2}\OperatorTok{{-}\textgreater{}}\NormalTok{data}\OperatorTok{.}\NormalTok{prim}\OperatorTok{.}\NormalTok{data}\OperatorTok{.}\NormalTok{nat}\OperatorTok{,} \OperatorTok{\&}\NormalTok{b2}\OperatorTok{{-}\textgreater{}}\NormalTok{data}\OperatorTok{.}\NormalTok{prim}\OperatorTok{.}\NormalTok{data}\OperatorTok{.}\NormalTok{nat}\OperatorTok{);} \OperatorTok{\}}
    \ControlFlowTok{return}\NormalTok{ stack}\OperatorTok{[}\NormalTok{sp }\OperatorTok{{-}} \DecValTok{1}\OperatorTok{];}
\OperatorTok{\}}

\NormalTok{Value }\OperatorTok{*}\NormalTok{dispatch}\OperatorTok{(}\NormalTok{Arena }\OperatorTok{*}\NormalTok{a}\OperatorTok{,} \DataTypeTok{int}\NormalTok{ proc\_id}\OperatorTok{,}\NormalTok{ Value }\OperatorTok{**}\NormalTok{env}\OperatorTok{,} \DataTypeTok{size\_t}\NormalTok{ env\_len}\OperatorTok{)} \OperatorTok{\{}
    \ControlFlowTok{switch} \OperatorTok{(}\NormalTok{proc\_id}\OperatorTok{)} \OperatorTok{\{}
        \ControlFlowTok{case} \DecValTok{0}\OperatorTok{:} \ControlFlowTok{return}\NormalTok{ proc\_0}\OperatorTok{(}\NormalTok{a}\OperatorTok{,}\NormalTok{ env}\OperatorTok{,}\NormalTok{ env\_len}\OperatorTok{);}
        \ControlFlowTok{case} \DecValTok{1}\OperatorTok{:} \ControlFlowTok{return}\NormalTok{ proc\_1}\OperatorTok{(}\NormalTok{a}\OperatorTok{,}\NormalTok{ env}\OperatorTok{,}\NormalTok{ env\_len}\OperatorTok{);}
        \ControlFlowTok{case} \DecValTok{2}\OperatorTok{:} \ControlFlowTok{return}\NormalTok{ proc\_2}\OperatorTok{(}\NormalTok{a}\OperatorTok{,}\NormalTok{ env}\OperatorTok{,}\NormalTok{ env\_len}\OperatorTok{);}
        \ControlFlowTok{case} \DecValTok{3}\OperatorTok{:} \ControlFlowTok{return}\NormalTok{ proc\_3}\OperatorTok{(}\NormalTok{a}\OperatorTok{,}\NormalTok{ env}\OperatorTok{,}\NormalTok{ env\_len}\OperatorTok{);}
        \ControlFlowTok{default}\OperatorTok{:}\NormalTok{ abort}\OperatorTok{();}
    \OperatorTok{\}}
\OperatorTok{\}}
\end{Highlighting}
\end{Shaded}

GCC inlines all four \texttt{proc\_N} functions into \texttt{dispatch}, and
inlines \texttt{op\_add\_nat} within the proc\_id=3 case. The fold loop calls
\texttt{call\_value} (external symbol -- closure dispatch), but the addition hot
path is branch-free for small nats. \texttt{op\_add\_nat} uses a two-branch
BigNat: values fitting in \texttt{uint64\_t} take a fast
\texttt{\_\_builtin\_add\_overflow} path, falling through to GMP only on
overflow.

\subsection{Why compile to C rather than LLVM or direct machine
code}\label{why-compile-to-c-rather-than-llvm-or-direct-machine-code}

\begin{itemize}
\tightlist
\item
  GCC's optimizer is free and excellent -- we get constant folding, inlining,
  register allocation, and instruction scheduling without implementing any of it
\item
  The generated C is trivial to read and debug
\item
  Cross-platform by construction (GCC targets everything)
\item
  The compilation overhead is large (\textasciitilde600ms-3s for GCC) but paid
  once; the \texttt{.so} is cached on disk keyed by a hash of the bytecode
\end{itemize}

\subsection{Arena allocation}\label{arena-allocation}

The native backend uses arena allocation: all Values created during one
execution are bump-allocated from a single arena, freed in one shot when the
call returns. This eliminates per-value malloc/free overhead and avoids the need
for a garbage collector.

\begin{Shaded}
\begin{Highlighting}[]
\CommentTok{// All allocation during execution goes through the arena}
\NormalTok{Value }\OperatorTok{*}\NormalTok{value\_alloc}\OperatorTok{(}\NormalTok{Arena }\OperatorTok{*}\NormalTok{a}\OperatorTok{);}
\DataTypeTok{void}  \OperatorTok{*}\NormalTok{arena\_alloc}\OperatorTok{(}\NormalTok{Arena }\OperatorTok{*}\NormalTok{a}\OperatorTok{,} \DataTypeTok{size\_t}\NormalTok{ bytes}\OperatorTok{);}
\CommentTok{// At call boundary: create arena, execute, extract result, free arena}
\end{Highlighting}
\end{Shaded}

\section{Design decisions}\label{design-decisions}

\textbf{Flat closures over linked environments}: every \texttt{CLOSURE} captures
the full current env as a vector. Wasteful in space (captures variables that may
never be used) but simple and cache-friendly. No linked-list traversal on
variable access, just \texttt{env{[}idx{]}}. The waste is bounded because Dhall
functions tend to be small.

\textbf{No jumps, only banks}: makes the instruction sequence trivially
serializable, trivially parallelizable per-bank, and trivially translatable to C
functions. The cost is extra procedure-call overhead for if/else, but Dhall
programs are not control-flow-heavy.

\textbf{Type erasure at compile time}: record field names, union constructor
names, and type arguments are all resolved to positional indices by the
compiler. The bytecode and interpreter never see strings for field access. This
is where the compile-time type information pays off -- we can use the Dhall type
checker as a compilation oracle.

\textbf{COMBINE/PREFER as pre-computed merge plans}: rather than interpreting
record merges dynamically (which requires knowing field names), the compiler
pre-computes a vector of ``take field idx from left/right'' operations. The
interpreter just follows the plan. \texttt{COMBINE} handles recursive merging by
encoding the full path as a \texttt{NESeq\ Word16} --
e.g.~\texttt{FromLeft\ (0\ :\textless{}\textbar{}\textbar{}\ {[}2{]})} means
``output field comes from left operand, field 0, then sub-field 2.''

\textbf{Builtins as partially-applied values}: \texttt{PUSH\_BUILTIN} pushes a
\texttt{VBuiltin} with an empty arg list. Each \texttt{CALL} adds one arg. When
the arg count hits the builtin's arity, it fires. This means builtins don't need
special syntax -- they're just values that happen to execute when saturated.
Partial application is free (it's just accumulating into a list).

\textbf{PType for de Bruijn alignment}: Dhall's type-level lambdas mean that
\texttt{\textbackslash{}(a\ :\ Type)\ -\textgreater{}\ \textbackslash{}(x\ :\ a)\ -\textgreater{}\ x}
has two parameters. At runtime, \texttt{a} is never inspected, but we still need
it to occupy env{[}0{]} so that \texttt{x} is at env{[}1{]} (since the compiler
uses de Bruijn indices). \texttt{PType} is the placeholder.

\textbf{Let as immediate closure application}: this might seem wasteful
(allocating a closure just to immediately call it), but it means the compiler
doesn't need a separate mechanism for let-bindings. The interpreter's
\texttt{CALL} path handles everything uniformly. In practice the overhead is
negligible because the ``closure'' is just the current env pointer (already
allocated).

\section{Signoff}\label{signoff}

Hi, thanks for reading! You can reach me via email at
\href{mailto:justin@jle.im}{\nolinkurl{justin@jle.im}}, or at twitter at
\href{https://twitter.com/mstk}{@mstk}! This post and all others are published
under the \href{https://creativecommons.org/licenses/by-nc-nd/3.0/}{CC-BY-NC-ND
3.0} license. Corrections and edits via pull request are welcome and encouraged
at \href{https://github.com/mstksg/inCode}{the source repository}.

If you feel inclined, or this post was particularly helpful for you, why not
consider \href{https://www.patreon.com/justinle/overview}{supporting me on
Patreon}, or a \href{bitcoin:3D7rmAYgbDnp4gp4rf22THsGt74fNucPDU}{BTC donation}?
:)

\textbackslash end\{document\}

\end{document}
